\documentclass[11pt,letterpaper]{article}
\usepackage[T1]{fontenc}
\usepackage[utf8]{inputenc}
\usepackage{lmodern}
\usepackage[margin=1in,headheight=15pt]{geometry}
\usepackage{amsmath,amssymb,amsthm,mathtools,mathrsfs}
\usepackage{microtype,booktabs,longtable,array,enumitem}
\usepackage[dvipsnames]{xcolor}
\usepackage{listings,fancyhdr}
\usepackage[colorlinks=true,linkcolor=MidnightBlue,citecolor=OliveGreen,urlcolor=MidnightBlue]{hyperref}
\usepackage[nameinlink,noabbrev]{cleveref}
\hypersetup{pdftitle={Surreal and Surcomplex Numbers in Symbolic Computer Algebra},pdfauthor={Research and implementation study},pdfsubject={Exact representations, effective subfields, Hahn series, transseries, and Wolfram Language architecture}}
\pagestyle{fancy}\fancyhf{}
\fancyhead[L]{\small Surreal and surcomplex computer algebra}
\fancyhead[R]{\small Exactness, algorithms, and semantics}
\fancyfoot[C]{\thepage}
\setlength{\parskip}{0.3em}
\setlength{\emergencystretch}{2em}
\setlist{itemsep=0.2em,topsep=0.35em}
\setcounter{tocdepth}{2}
\newtheorem{theorem}{Theorem}[section]
\newtheorem{proposition}[theorem]{Proposition}
\newtheorem{lemma}[theorem]{Lemma}
\newtheorem{corollary}[theorem]{Corollary}
\theoremstyle{definition}
\newtheorem{definition}[theorem]{Definition}
\newtheorem{example}[theorem]{Example}
\newtheorem{remark}[theorem]{Remark}
\newcommand{\No}{\mathbf{No}}
\newcommand{\SC}{\mathbf{No}[i]}
\newcommand{\R}{\mathbb R}
\newcommand{\C}{\mathbb C}
\newcommand{\Q}{\mathbb Q}
\newcommand{\Z}{\mathbb Z}
\newcommand{\N}{\mathbb N}
\newcommand{\RA}{\mathbb R_{\mathrm{alg}}}
\newcommand{\OO}{\mathcal O}
\newcommand{\mm}{\mathfrak m}
\newcommand{\HH}{\mathcal H}
\newcommand{\st}{\operatorname{st}}
\newcommand{\fin}{\operatorname{fin}}
\newcommand{\supp}{\operatorname{supp}}
\newcommand{\lc}{\operatorname{lc}}
\newcommand{\sgn}{\operatorname{sgn}}
\newcommand{\cis}{\operatorname{cis}}
\newcommand{\Arg}{\operatorname{Arg}}
\newcommand{\Sin}{\operatorname{Sin}}
\newcommand{\Cos}{\operatorname{Cos}}
\newcommand{\ExpNo}{\operatorname{Exp}_{\No}}
\newcommand{\LogNo}{\operatorname{Log}_{\No}}
\newcommand{\Expi}{\operatorname{Exp}_{\mathrm{inf}}}
\newcommand{\Logi}{\operatorname{Log}_{\mathrm{inf}}}
\newcommand{\re}{\operatorname{Re}}
\newcommand{\im}{\operatorname{Im}}
\newcommand{\Res}{\operatorname{Res}}
\newcommand{\Ov}[1]{O_v(#1)}
\newcommand{\sh}{^{\sharp}}
\newcommand{\HInt}{\mathop{\int^{\!H}}}
\newcommand{\code}[1]{\texttt{\detokenize{#1}}}
\newcolumntype{L}[1]{>{\raggedright\arraybackslash}p{#1}}
\lstdefinestyle{wl}{basicstyle=\small\ttfamily,breaklines=true,columns=fullflexible,keepspaces=true,showstringspaces=false,frame=single,framerule=0.3pt,rulecolor=\color{gray!45},backgroundcolor=\color{gray!3},xleftmargin=0.4em,xrightmargin=0.4em,aboveskip=0.7em,belowskip=0.7em}
\lstset{style=wl}
\numberwithin{equation}{section}
\begin{document}
\begin{titlepage}
\centering
\vspace*{0.6cm}
{\Large\scshape Mathematical semantics and implementable algorithms\par}
\vspace{0.85cm}
{\Huge\bfseries Surreal and Surcomplex Numbers\par}
\vspace{0.35cm}
{\LARGE\bfseries in Symbolic Computer Algebra\par}
\vspace{0.8cm}
{\Large Exact Representations, Effective Subfields,\par Hahn Series, Transseries, and a Wolfram Language Design\par}
\vspace{0.8cm}
{\large September 21, 2026\par}
\vspace{0.6cm}
\begin{minipage}{0.94\textwidth}
\begin{abstract}
A symbolic computer algebra system can support substantial and useful parts of the surreal field and its complexification without attempting to store arbitrary surreal numbers. The decisive distinction is between an exact expression that denotes a number, a representation with effective coefficient access, and a representation supporting terminating equality, order, and normalization algorithms. These are different capabilities.

We develop a layered design grounded in the supplied manuscripts on surcomplex analysis and trigonometry. A finite-rank rational-function field, with algebraic real coefficients and a certified order on monomials, provides a decidable algebraic core; its effective real closure and complexification support exact polynomial algebra. Certified Hahn-series constructors add infinitesimal Taylor evaluation, lazy reciprocal and algebraic expansions, and support-aware precision. Exponential--logarithmic expressions and transseries enlarge the range further, while keeping their interpretation and normalization obligations explicit. Canonical finite-angle and strip trigonometry, branch-sensitive logarithms, implicit functions, coherent contour functionals, and scalar derivations receive separate interfaces.

We prove elementary implementation-boundary results, including an undecidability obstruction for unrestricted computable coefficient streams and an obstruction to finite-list truncation in higher-rank value groups. Detailed algorithms, domain tables, worked examples, a staged implementation plan, and a small Wolfram Language reference kernel make the proposal concrete. The kernel implements only its advertised rational subfield; it is not a universal surreal CAS. No claim of complete symbolic normalization, unrestricted convergence, or a canonical infinite circular phase is made.
\end{abstract}
\end{minipage}
\vfill
\begin{minipage}{0.94\textwidth}\small
\textbf{Provenance and status.} The two supplied manuscripts determine the analytic conventions. Published background, manuscript-specific results, deductions proved here, and proposed software interfaces are distinguished in the text. The included finite symbolic checks are not a formal verification of the theory. External documentation and research records were consulted on September 21, 2026.
\end{minipage}
\end{titlepage}
\begingroup\footnotesize\setlength{\parskip}{0pt}\tableofcontents\endgroup
\clearpage

\section{What a surreal CAS should promise}\label{sec:goals}
A useful first target is not ``all surreal numbers in a new numeric head.'' It is an extensible family of exact domains with explicit guarantees. For example, the following calculations should be straightforward in an initial implementation:
\begin{align*}
 (\omega+1)^2-\omega^2-2\omega&=1,\\
 (\omega+1)^{-1}&=\frac{t}{1+t},\qquad t=\omega^{-1},\\
 \sqrt{\omega^2+1}&=\omega\sqrt{1+t^2},\\
 \arctan\omega&=\frac\pi2-\arctan t.
\end{align*}
The second equality is already an exact answer. Expanding it into infinitely many monomials is optional. The third requires an algebraic extension, and the fourth requires the canonical finite-angle analytic layer. A system should be able to explain these changes of domain.

Three questions guide the design. What mathematical object does the expression denote? What operations can the representation carry out with guaranteed termination? What information is retained when the system returns a finite approximation? Answering the first does not settle the other two.

\subsection{The supplied manuscripts as the baseline}
The analysis manuscript \cite{SourceAnalysis} distinguishes fine-local formal analytic germs, Hahn-coherent functions on thickenings of ordinary complex domains, and a much stronger all-scale coherence condition. It defines infinite evaluation using strong Hahn summability, not ordinary convergence of partial sums. The trigonometry manuscript \cite{SourceTrig} supplies canonical sine and cosine on finite real surreals, polar decomposition at arbitrary scales, inverse trigonometry, canonical strip functions, and a classification of additional infinite-phase choices.

We retain those distinctions. In particular, the coefficientwise contour integral below is the manuscript's functional, not a newly asserted full-class Riemann integral. Its all-scale rigidity theorems are results of its specified category, not a prohibition on transcendental surreal functions. The two other archives mentioned inside the trigonometry manuscript were not supplied for this study and are not treated as independently reviewed sources.

The mathematical foundation includes the normal-form and ordered exponential structures developed in the literature \cite{Conway,Gonshor,vdDE}, together with Hahn-field algebra \cite{Poonen}. The effective-domain construction, software contracts, and implementation recommendations below are the present article's synthesis. The small reference program demonstrates only a selected part of that design.

\subsection{A concise answer to the representation question}
There is no largest useful finitely representable subset independent of a chosen language: a language can always acquire an additional exact constant or certified constructor. There are, however, natural tiers. Finite monomial sums and rational functions admit finite algebraic storage. Algebraic extensions admit finite defining equations and root selectors. Analytic values at infinitesimal perturbations admit exact germ descriptions and lazy coefficients. Exponential--logarithmic and transserial expressions admit exact structured descriptions with more demanding comparison and support problems. Arbitrary surreal numbers, arbitrary coefficient streams, and arbitrary set-sized supports do not admit a universal finite effective representation.

The right product is therefore a \emph{decidable core plus explicitly partial extensions}. It should never silently replace an unsupported value by zero, ordinary infinity, a large real approximation, or a convention-dependent phase.

\section{Mathematical objects, and distinctions that affect code}\label{sec:semantics}
\subsection{Normal forms and workspaces}
We use the convention of both supplied manuscripts:
\begin{equation}\label{eq:normal}
 t^\gamma:=\omega^{-\gamma},\qquad
 a=\sum_{\gamma\in S}a_\gamma t^\gamma,
 \qquad S\subseteq\Gamma\text{ well ordered}.
\end{equation}
Here the displayed power of $\omega$ is the \emph{Conway monomial map}. For real surreal scalars, $a_\gamma\in\R$; for surcomplex scalars, $a_\gamma\in\C$. The ordered additive group $\Gamma$ is a set-sized subgroup of $\No$. Every collection of coefficients, every support, and every family being summed in this article is a set. The full $\No$ is a proper class, but every individual expression is finite data interpreted in suitable set-sized workspaces.

If $a\ne0$, define
\[
 v(a)=\min\supp(a),\qquad \lc(a)=a_{v(a)},\qquad v(0)=+\infty.
\]
Then $v(ab)=v(a)+v(b)$ and $v(a+b)\ge\min(v(a),v(b))$. For real $a$, its sign is the sign of its leading coefficient. For a divisible $\Gamma$, $\C((t^\Gamma))$ is algebraically closed; the corresponding real Hahn field is real closed. The algebraic-closure input is a theorem about the complete mathematical field, not a terminating algorithm for arbitrary coded series \cite[Corollary 4]{Poonen}.

\subsection{Three kinds of power and two kinds of infinity}
A software interface should distinguish at least
\[
 \code{ConwayMonomial[a]},\qquad
 \code{SurrealPower[x,y]},\qquad
 \code{OrdinalPower[a,b]}.
\]
For positive surreal $x$, the second may mean $\ExpNo(y\LogNo x)$. The first is a normal-form constructor, and the third is ordinal exponentiation. Identities relating them must be proved on their particular domains, rather than inferred from identical typography. Likewise, ordinary ordinal addition is not surreal field addition: $1+_{\mathrm{ord}}\omega=\omega$, whereas $1+\omega=\omega+1$ in the field. Conway's embedding relates field arithmetic on ordinals to natural ordinal arithmetic, not to an unrestricted replacement of ordinal operations by field operations \cite{Conway}.

A genuine infinite field element is also not a projective point at infinity. In particular,
\[
 \omega-\omega=0,\qquad \omega/\omega=1,\qquad \omega^{-1}>0.
\]
A head used for a pole or an infinite limit cannot supply those field laws. Nor is $t$ a nilpotent dual number: $t^n\ne0$ for every ordinary $n$. Automatic-differentiation jets and surreal infinitesimals have different multiplication semantics.

\subsection{Standard part is partial; finite part is not multiplicative}
Let $\OO$ denote the finite surreal or surcomplex elements and $\mm$ the infinitesimals. Every $z\in\OO$ has a unique decomposition
\[
 z=c+\eta,\qquad c\in\R\text{ or }\C,\quad\eta\in\mm.
\]
The standard part $\st(z)=c$ is a ring homomorphism \emph{on this finite subring}. The finite-part projection $\fin$ used in the trigonometry manuscript instead removes all negative-exponent terms from any normal form. It is additive, but
\[
 \fin(\omega\omega^{-1})=1\ne
 \fin(\omega)\fin(\omega^{-1})=0.
\]
Extracting the coefficient at exponent zero from arbitrary infinite elements is not a ring homomorphism either. These operations require different names and rewrite rules. Even when a standard part exists mathematically, extracting it from an arbitrary program-defined series may not be computable.

\subsection{A type is more than a coefficient container}
An exact object should carry or reference its coefficient domain, exponent domain and order, embedding into surreal numbers, analytic interpretation, and any branch or derivation convention. A formal generator called $s$ is not enough: $s=t^2$ and $s\prec t^n$ for every ordinary $n>0$ give different fields and different equality tests. Domain identifiers should be structural and serializable, not dependent on session-local assumptions that can later change.

\section{Four levels of exactness and two unavoidable limits}\label{sec:exactness}
\subsection{Denotation, effective access, decision, and approximation}
\begin{definition}[Representation capabilities]
A finite description is \emph{denotationally exact} when its interpretation uniquely specifies a mathematical value. It is \emph{coefficient-effective} relative to an indexed support description when requested coefficients can be computed by terminating procedures. It is \emph{decision-effective} for a specified collection of operations when those operations, such as equality or real order, always terminate correctly. A \emph{certified truncation} describes a set of possible exact values by a remainder condition; it is not itself an exact value unless the remainder is certified zero.
\end{definition}
For instance, an exact root specified by a polynomial and a valid root selector may be easier to compare than an unrestricted lazy power series, although the latter offers direct coefficient access. A name such as $\ExpNo(\omega+\omega^{-1})$ can be unambiguous without its entire normal form being available. A symbolic integral can name an exact value without an algorithm simplifying that value into a preferred grammar.

The capabilities should be queried independently. A node may support exact addition and multiplication by building expression graphs but offer only a partial equality procedure. In that layer, a time-limited unsuccessful comparison is \code{Unknown}, not \code{False}. Mathematical equality, syntactic identity, and equality modulo a remainder ideal must also be distinct operations.

\subsection{There cannot be finite names for every surreal}
\begin{proposition}[Finite-language cardinality limit]\label{prop:countability}
Fix a countable collection of symbols, each itself specified by finite data. The values denoted by its finite expressions form at most a countable set. Hence such a language cannot name every real number, and a fortiori cannot name every surreal number.
\end{proposition}
\begin{proof}
There are countably many finite strings over a countable alphabet. Taking the image of the partial interpretation map cannot increase cardinality. The ordinary reals are already uncountable.
\end{proof}
An oracle or an opaque atom can be allowed to stand for an arbitrary real or surreal parameter, but the missing information has then been moved into an oracle or an external assumption. It has not been encoded in ordinary finite memory. Similarly, ``every normal form has set-sized support'' is not a computability statement. A set may be uncountable, noncomputable, or lack any finite structural description.

\subsection{Even total computable coefficient streams defeat universal equality}
\begin{theorem}[A zero-test obstruction]\label{thm:halting}
There is no algorithm which, for every total computable sequence $(a_n)_{n\ge0}$ with $a_n\in\{0,1\}$, decides whether the Hahn series $\sum_{n\ge0}a_nt^n$ is zero. Consequently no universal terminating equality or three-way sign procedure exists for this unrestricted representation class.
\end{theorem}
\begin{proof}
For a given machine $M$ and its input, let $a_n=1$ if $M$ has halted within $n$ steps, and $a_n=0$ otherwise. Computing $a_n$ requires only bounded simulation, so this is a total computable sequence. Its support is contained in $\N$, making the displayed series a valid Hahn sum at $t=\omega^{-1}$. It is zero exactly when $M$ never halts. A total zero test would therefore decide the halting problem. All coefficients are nonnegative, so a sign procedure distinguishing zero from positive would do the same. Equality reduces to zero testing by subtraction.
\end{proof}
In this reduction every individual value is either zero or a rational function $t^N/(1-t)$ for some halting time $N$. Thus even values lying in a decidable rational field can become undecidable when presented by unrestricted coefficient programs. Representation, not merely membership in an abstract field, determines the algorithmic guarantee.

This obstruction requires neither arbitrary real coefficients nor complicated surreal exponents. A certified support container $\N$ is available from the outset. The problem is locating the first nonzero coefficient or proving that none exists. A lazy series implementation can still be useful, but must admit nontermination, resource-bounded \code{Unknown}, or restricted input grammars with extra zero certificates. Sage's lazy-series documentation makes the related practical distinction between coefficient generation and equality decidable only in special cases \cite{SageLazy}.

The result does not prevent a zero test for rational generating functions, algebraic series with finite defining data, or other specifically certified families. It says that the promise must be attached to the representation class rather than to the word ``exact.''

\section{A taxonomy of exact representations}\label{sec:representations}
\subsection{Cuts and sign sequences}
Conway cuts and sign sequences are closest to the foundational construction. Finite sign strings represent the short surreal numbers, namely dyadic rationals. They are valuable for teaching, checking small recursive definitions, and recording simplicity or birthday information. They are a poor universal arithmetic storage format: the ordinary number $1/3$ has a short rational expression but is not a finite-birthday surreal.

Compressed ordinal-indexed sign blocks can name selected transfinite objects. Their usefulness depends on algorithms for their ordinal notations and block operations. A finite recursive cut description may also name some infinite numbers, but comparing arbitrary presented cuts can require universal comparisons with infinitely many options. Neither birthday nor simplicity is determined by the field structure alone. A Hahn-field backend that knows addition and order should not pretend automatically to provide shortest sign expansions or the simplest number between arbitrary cuts.

\subsection{Finite and hereditary normal forms}
A sparse finite sum $\sum_{j=1}^m c_jt^{\gamma_j}$ has a finite representation once coefficients and exponents do. Addition collects equal exponents, multiplication adds exponents, and ordering uses the first surviving coefficient. Recursive finite normal forms for the exponents extend the hierarchy: a coefficient/exponent tree can describe nested scales without materializing a sign sequence.

This representation is generally a ring, not a field. For example,
\[
 (1-t)^{-1}=\sum_{n\ge0}t^n
\]
is not a finite-support normal form. An implementation must either move to rational-expression storage, create a lazy infinite sum with a support certificate, or return an explicit unsupported-domain result. Declaring the inverse to have a finite normal form would be incorrect.

\subsection{Rational expressions and algebraic root data}
Fractions of finite monomial polynomials are especially effective. Their arithmetic can often be performed without any series expansion. Exact algebraic extensions then use polynomial equations and root selectors. A root selector might consist of a Thom encoding, a certified infinitesimal expansion distinguishing a branch, or a tower of ordered algebraic extensions. A numeric approximation to an ordinary real root is not a sufficient selector for roots separated only by infinitesimals.

The output of an algebraic operation need not belong to the input's finite-support grammar. It can nevertheless have a finite exact implicit representation. Thus $\sqrt{1+t}$ is most naturally stored as the positive root of $X^2-1-t$, with its binomial expansion an auxiliary view rather than its definition.

\subsection{Puiseux, grid-based, and general Hahn data}
In a one-parameter characteristic-zero setting, algebraic branches can be described by Puiseux expansions after ramification. The \emph{field of all Puiseux series} still permits arbitrary infinite coefficient sequences; it should not be identified with the finitely described algebraic functions. Finite defining equations plus branch data give an effective class, whereas an arbitrary coefficient stream leads back to \cref{thm:halting}.

A grid-supported Hahn expression has support contained in
\begin{equation}\label{eq:grid}
 \gamma_0+\N\lambda_1+\cdots+\N\lambda_r,
 \qquad \lambda_j>0.
\end{equation}
This finite container makes many closure arguments explicit. Different tuples may represent the same exponent; coefficients must be combined through finite convolution, not treated as distinct monomials. Grid-based methods have a substantial existing computational theory \cite{vdH}. They do not make every coefficient sequence decidable, nor do they imply that every valuation cutoff contains only finitely many terms.

General Hahn supports can be well ordered without lying in such a finite grid. Even a countable support of order type $\omega+1$ cannot be completely traversed by a monotonically increasing ordinary stream: after its first $\omega$ terms, no finite step reaches the final term. Such supports require nested blocks, symbolic support sets, indexed access, or exact structural expressions. ``Lazy'' must not be used as a synonym for ``an ordinary iterator suffices.''

\subsection{Expression graphs, transseries, and hyperseries}
A directed acyclic graph built from exact constants, field operations, algebraic roots, selected exponentials and logarithms, and certified sums is a flexible denotational language. Its finite size does not imply a finite normal form. Common-subexpression sharing is particularly valuable for algebraic towers and repeated local expansions.

Transseries offer a structured enlargement for exponential and logarithmic scales and for composition at infinity. Published work interprets LE-, EL-, and omega-series as functions on positive infinite surreals, with composition and differentiation in specified classes \cite{BMFunctions}. For implementation, select finite or effectively generated fragments with explicit coefficient and support algorithms; the full mathematical classes are not all finitely described.

The hyperseries representation theorem of Bagayoko and van der Hoeven provides a much broader mathematical correspondence between surreal values and growth-rate expressions \cite{Hyperseries}. It is not a finite-string encoding of every surreal. The countability argument still applies, and a representation theorem with infinitary data does not by itself supply normalization, equality, or a storage format.

\begin{center}\small
\begin{longtable}{@{}L{0.22\textwidth}L{0.34\textwidth}L{0.35\textwidth}@{}}
\toprule
Representation & Useful exact content & Principal implementation boundary\\\midrule\endhead
Finite sign data & Dyadics and small foundational examples & Does not include all rationals; recursive cost and birthday issues.\\[0.35em]
Finite normal forms & Sparse monomial rings and nested finite scales & Usually not closed under reciprocal.\\[0.35em]
Ordered rational functions & Fractions in finitely many certified monomials & Algebraic independence and monomial order must be fixed.\\[0.35em]
Ordered algebraic data & Real closure and complex polynomial roots & Needs non-Archimedean root selectors, not floating-point isolation.\\[0.35em]
Certified series grammars & Exact infinite sums and local Taylor values & Support validity and effective zero/leading-term access are separate.\\[0.35em]
Exp--log/transserial graphs & Infinite exponential scales, asymptotic composition & Partial normalization; domain- and interpretation-specific operations.\\[0.35em]
General oracular Hahn data & Semantic models and conditional symbolic calculations & No universal effective equality or finite cutoff materialization.\\
\bottomrule
\end{longtable}
\end{center}

\section{A concrete decidable algebraic core}\label{sec:core}
\subsection{Independent monomials with a computable order}
Fix an ordinary positive integer $d$ and define
\begin{equation}\label{eq:generators}
 \gamma_j=\omega^{j-1},\qquad t_j=t^{\gamma_j}
       =\omega^{-\omega^{j-1}},\qquad 1\le j\le d.
\end{equation}
These exponents are $\Q$-linearly independent. The sign of $\sum_jq_j\gamma_j$ is the sign of its last nonzero coefficient. Consequently
\[
 0<t_{j+1}<t_j^n\qquad(n\in\N_{>0}),
\]
where $n$ is always an ordinary finite integer. The exponent vector $(q_1,\ldots,q_d)$ is ordered by comparing its highest-index differing coordinate first. More general finitely generated ordered exponent groups are possible, but they need an equally explicit equality and order algorithm.

Set
\begin{equation}\label{eq:Fd}
 F_d=\RA(t_1,\ldots,t_d)\subset\No.
\end{equation}
The coefficient field $\RA$ consists of real algebraic numbers, with exact arithmetic and order. An initial program may use $\Q$ instead, as our reference kernel does.

\begin{proposition}[Faithfulness and rational-field algorithms]\label{prop:core}
The substitution in \eqref{eq:generators} embeds the ordinary rational-function field in $d$ algebraically independent variables into $\No$. Equality, arithmetic, real order, valuation, leading coefficient, and standard part of finite elements of $F_d$ are computable by terminating finite algorithms.
\end{proposition}
\begin{proof}
Distinct integer exponent vectors give distinct $\sum n_j\gamma_j$ by independence. Thus a nonzero polynomial maps to a nonempty finite normal form, so cannot vanish. The substitution is injective and extends to fractions.

For a nonzero polynomial $P$, collect its monomials and choose the least exponent vector $\alpha(P)$ in the specified order; let $c(P)$ be that coefficient. For $f=P/Q\ne0$,
\begin{equation}\label{eq:leadingrational}
 v(f)=\alpha(P)-\alpha(Q),\quad
 \lc(f)=c(P)/c(Q),\quad
 \sgn(f)=\sgn(c(P)c(Q)).
\end{equation}
Polynomial arithmetic and coefficient comparisons are finite. Equality of two fractions is decided by the zero polynomial test on $PQ'-P'Q$. If $v(f)>0$, then $\st(f)=0$; if $v(f)=0$, then $\st(f)=\lc(f)$; if $v(f)<0$, then $f$ is infinite and the finite standard-part operation is inapplicable. The zero case is handled separately.
\end{proof}
Cancellation before leading-term extraction is essential. For example, the numerator $(1+t_1)t_2-t_2$ has leading exponent $(1,1)$, not $(0,1)$. A heuristic inspection of the uncollected expression would be wrong.

\subsection{Effective real closure and complexification}
Let $R_d$ be the real closure of this ordered field, realized inside $\No$, and put $C_d=R_d[i]$. These are much stronger targets than finite monomial sums: $R_d$ is closed under positive square roots and roots of odd-degree polynomials, while $C_d$ is algebraically closed.

An exact real-root object can store a polynomial over a previously constructed effective ordered field and a Thom encoding: the signs of its successive derivatives at the selected real root. The selector distinguishes a root algebraically, even when ordinary numerical information cannot. A finite tower of such objects is an alternative to reducing every operation immediately to a single primitive element.

\begin{theorem}[An implementable real-closed algebraic layer]\label{thm:realclosure}
There is a finite-data representation of $R_d$, and hence of $C_d$, supporting terminating exact field arithmetic, equality, real order on $R_d$, and all roots of ordinary finite-degree univariate polynomials. Semialgebraic sentences over these effective coefficient fields can be decided. No Hahn expansion of each algebraic element is required for these guarantees.
\end{theorem}
\begin{proof}[Algorithmic justification]
Use finite polynomial equations with unique ordered-root selectors for algebraic elements. Addition, multiplication, and inversion introduce finite systems expressing the operation and the selectors of its arguments. Elimination produces equations for the result; sign determination chooses the intended real root. Equality and order are likewise finite formulas involving the defining equations and selectors. Real-closed-field quantifier elimination reduces each such formula to finitely many polynomial sign tests over the effective base field. Those tests terminate by \cref{prop:core}. Algebraic sign-determination and Thom-encoding elimination algorithms are available independently of Archimedean topology \cite{PerrucciRoy}.

Pairs of these real elements represent $C_d$. A polynomial equation in one complex variable becomes a pair of real polynomial equations; its finite root set and coordinate selectors can be handled by the same finite algebraic machinery. A positive-dimensional multivariable solution set requires a symbolic set description, not a finite list of points. This constructs effective operations on the specified real closure and its complexification, rather than on arbitrary externally supplied surreal oracles.
\end{proof}
The theorem establishes implementability, not acceptable complexity in every case. Elimination can be expensive and degree growth severe. In practice, maintain extension towers and specialized univariate routines, and use general elimination as a completeness fallback.

\subsection{Why ordinary isolating intervals can be inadequate}
Let $F=\RA(t)$ and let $t>0$ be infinitesimal. The two real-closure elements
\[
 \alpha=\sqrt t,\qquad \beta=\sqrt t+t
\]
are distinct. There is no element of $F$ strictly between them. Indeed, a positive rational function in $t$ has an integer leading exponent. If that exponent is at least $1$, the value is smaller than $\sqrt t$; if it is at most $0$, the value is larger than $\sqrt t+t$. Thus a root-isolation strategy that insists on refining endpoints in $F$ cannot separate this pair.

This is an algorithm-design obstruction, not a failure of ordered algebraic root representation. Thom encodings or suitable algebraic extension endpoints resolve it. In particular, increasing the decimal precision of a real number cannot distinguish two roots with identical real standard part and infinitesimal separation.

\subsection{What transfer does, and does not, buy}
The ordered-field laws and real-closed algebraic theorems transfer to this core. Finite matrix identities, polynomial identities, and semialgebraic geometry can therefore be reused with a coefficient-domain adapter. The statement that $t$ is smaller than $1/n$ for every \emph{ordinary integer} $n$ is not a single ordinary ordered-field formula defining an infinitesimal in $\R$. Passing a few inequalities to a real solver is not an implementation of that condition.

Likewise, an elementary-extension theorem for a surreal exponential structure is a semantic theorem, not a universal decision algorithm for exponential expressions \cite{vdDE}. A CAS must not infer that its ordinary real simplifier now knows every surreal exponential identity.

\section{Hahn summability as a checked constructor}\label{sec:support}
\subsection{The two conditions are both necessary}
A family $(A_j)_{j\in J}$ is \emph{strongly summable} when its union of supports is well ordered and every exponent receives contributions from only finitely many family members. Both manuscripts use precisely this meaning. The sum is formed by finite coefficient addition at each exponent \cite{SourceAnalysis,SourceTrig}.

For well-ordered supports $S,T$, the support of a product is contained in $S+T$, with finitely many contributing pairs at each exponent. If $E$ is well ordered and strictly positive, the monoid $E^*$ of finite sums from $E$ is well ordered and each exponent has finitely many word decompositions, up to the usual finite permutations. This Hahn--Neumann calculus justifies reciprocal series, infinitesimal Taylor substitution, and compatible coefficient regrouping.

An implementation should construct these certificates compositionally. A finite support supplies a base certificate; product uses a Minkowski-sum constructor; inversion of a normalized infinitesimal perturbation uses a positive-monoid constructor; analytic substitution combines the outer support with such a monoid. A bare statement that a user-supplied program ``probably produces increasing exponents'' is not such a certificate.

\subsection{Examples a correct evaluator must distinguish}
For $t=\omega^{-1}$,
\begin{align*}
 \sum_{n\ge0}t^n&=\frac1{1-t}, &
 \sum_{n\ge0}n!t^n&\text{ is a valid Hahn sum},\\
 \sum_{n\ge0}t^{-n}&\text{ is not a Hahn sum}, &
 \sum_{n\ge0}\frac1{n!}&\text{ is not a strongly summable Hahn family}.
\end{align*}
The third support decreases indefinitely. In the fourth family, infinitely many summands use exponent zero. Its ordinary real sum is, of course, $e$, and the CAS may evaluate that \emph{ordinary coefficient sum} first and embed the result. That is a different operation, not a relaxation of finite coefficient multiplicity in arbitrary Hahn sums.

The factorial series has zero ordinary radius of convergence but positive, well-ordered infinitesimal support. Its exact surreal value does not automatically identify it with a particular analytic continuation or summation of an ordinary divergent function. A special-function or resurgent-extension provider must supply that extra interpretation.

\subsection{A support certificate is not yet an evaluation algorithm}
To compute a requested coefficient, a backend also needs an effective way to identify all contributing decompositions and to know when the list is complete. A theorem that each list is finite does not automatically provide a terminating search that detects its end. Useful constructors therefore carry finite dependency bounds or explicit decomposer procedures in addition to semantic well-ordering evidence.

Similarly, a transfinite coefficient recursion proves the existence of a complete Hahn object without being an ordinary finite loop through the entire support. The correct computational implementation either produces a finite implicit object, computes specified coefficients with certified finite dependencies, or works inside a restricted effective support grammar. It must not report that it has ``finished the transfinite recursion'' merely because it computed many initial coefficients.

\subsection{The manuscript's small-scale theorem and effectivity}
The analysis manuscript proves that an arbitrary formal series $\sum a_nX^n$ with set-many surcomplex coefficients can be evaluated after a sufficiently small rescaling. Its proof chooses a surreal exponent dominating a set of exponents from \emph{all} coefficients and shifts their supports into separated bands \cite[\S4, theorem ``Evaluation of arbitrary formal series'']{SourceAnalysis}.

This is a valuable semantic existence theorem, but not a generic finite procedure for discovering that bound from an arbitrary coefficient program. An effective interface should ask for a support-growth bound, a dominating-scale certificate, or a restricted coefficient grammar from which one can be derived. A constructed scale may also lie outside the user's current workspace. Domain enlargement must be explicit rather than hidden inside a purported fixed-field operation.

\section{Exact arithmetic, lazy expansion, and precision}\label{sec:arithmetic}
\subsection{Addition and multiplication}
For finite sparse normal forms, addition is a merge after exponent normalization; multiplication generates pairs, adds their exponents, multiplies their coefficients, and collects equal exponents. With $m,n$ nonzero input terms, direct multiplication creates at most $mn$ pairs. This is an arithmetic-operation count, not a bit-complexity estimate: exponent comparison and coefficient growth can dominate the actual cost.

For lazy Hahn data, the same coefficient equations are valid, but each operation needs its certified dependency enumerator. Cancellation may remove an arbitrarily long initial segment. A budget-limited leading-term query should return an unevaluated request or \code{Unknown} rather than guessing that the next inspected coefficient is the leader.

For rational-function storage, use exact polynomial numerator--denominator arithmetic and cancellation; there is usually no reason to expand the result merely to add it to another fraction. Maintain both the compact exact expression and any requested series view, so that an expansion cache never becomes the only surviving definition.

\subsection{Inversion and rational powers}
If leading data are available, write a nonzero element as
\begin{equation}\label{eq:unit}
 a=c\,t^\gamma(1+\eta),\qquad c\ne0,\quad v(\eta)>0.
\end{equation}
Then
\begin{equation}\label{eq:inverse}
 a^{-1}=c^{-1}t^{-\gamma}\sum_{n\ge0}(-\eta)^n.
\end{equation}
The exact reciprocal may be retained as a rational node; \eqref{eq:inverse} is the certified expansion rule. For positive real $a$ and an ordinary positive integer $q$,
\begin{equation}\label{eq:rootbinomial}
 a^{1/q}=c^{1/q}t^{\gamma/q}
       \sum_{n\ge0}\binom{1/q}{n}\eta^n,
\end{equation}
with the positive coefficient root. Adjoining $\gamma/q$ may enlarge the exponent group, and adjoining $c^{1/q}$ may enlarge the coefficient field. Neither enlargement should be concealed.

For complex $a$, a chosen $q$th root of $c$ gives one branch and the others are obtained by multiplying by ordinary $q$th roots of unity. A branch selector remains necessary even though algebraic closedness supplies all roots. Formulas such as $(ab)^{1/q}=a^{1/q}b^{1/q}$ must not be used with unrelated complex root choices.

For a normalized one-parameter series $A(t)=1+\sum_{n\ge1}a_nt^n$, the reciprocal coefficients are obtained without symbolic differentiation:
\[
 b_0=1,\qquad b_n=-\sum_{k=1}^n a_kb_{n-k}.
\]
Each requested $b_n$ has finitely many dependencies. In a general Hahn group, the corresponding recursion is indexed by actual exponents; a positive-support certificate supplies the well-founded relation, but a computational decomposer is still required.

\subsection{A precise remainder convention}
We use
\[
 x=p+\Ov{\alpha}\quad\Longleftrightarrow\quad v(x-p)\ge\alpha,
\]
with $v(0)=+\infty$. This is an exact enclosure statement in the chosen valued workspace. It must not be confused with ordinary real asymptotic big-$O$ or with a statement that the remainder vanishes.

\begin{proposition}[Basic valuation-precision propagation]\label{prop:precision}
If $x=p+\Ov{\alpha}$ and $y=q+\Ov{\beta}$, then
\begin{align}
 x+y&=p+q+\Ov{\min(\alpha,\beta)},\\
 xy&=pq+\Ov{\min(\alpha+v(q),\,\beta+v(p),\,\alpha+\beta)}.
 \label{eq:productprecision}
\end{align}
If $p\ne0$ and $\alpha>v(p)$, then $x\ne0$ and
\begin{equation}\label{eq:inverseprecision}
 x^{-1}=p^{-1}+\Ov{\alpha-2v(p)}.
\end{equation}
\end{proposition}
\begin{proof}
Write $x=p+e$, $y=q+f$ and expand. The error in the product is $pf+qe+ef$, giving \eqref{eq:productprecision}. Under $v(e)>v(p)$, the leading term of $x$ equals that of $p$, so $v(x)=v(p)$. The identity $x^{-1}-p^{-1}=-e/(xp)$ gives \eqref{eq:inverseprecision}.
\end{proof}
A center polynomial that cancels to zero does not prove the exact answer zero. For instance, subtracting two identical truncations of independently represented numbers leaves a remainder enclosure unless exact equality of the underlying objects has been proved. Conversely, subtracting an object from itself should use dependency information and return exact zero rather than two independent error bounds.

Valuation precision and precision of ordinary real coefficients are independent. An interval coefficient enclosure containing zero does not determine a leading sign. A serious system should record both the exponent cutoff and the coefficient-error model, and refine whichever prevents the requested conclusion.

\subsection{Higher rank defeats finite-list cutoffs}
\begin{theorem}[A finite truncation request can have infinite output]\label{thm:cutoff}
In an ordered group containing positive $\lambda,\mu$ with $n\lambda<\mu$ for every ordinary $n$, a valid grid-supported Hahn series can have infinitely many terms below a fixed valuation cutoff. Consequently no general operation can always return all terms below a cutoff as a finite list.
\end{theorem}
\begin{proof}
Consider
\[
 A=\sum_{n\ge0}t^{n\lambda}+t^\mu
   =\frac1{1-t^\lambda}+t^\mu.
\]
The support $\{n\lambda:n\ge0\}\cup\{\mu\}$ is well ordered and lies in the grid generated by $\lambda,\mu$. All the first summands have exponents below $\mu$. Thus the truncation below $\mu$ has infinitely many nonzero terms. A finite list cannot represent that support completely.
\end{proof}
There is no obstruction to an \emph{exact compressed answer}: the truncation is $(1-t^\lambda)^{-1}$. This suggests three distinct interfaces: a finite number of leading terms; a valuation-cutoff object whose representation may remain compressed; and a finite rectangular or total-degree jet in selected generators. These are different requests.

After $N$ geometric terms, the remainder has valuation $(N+1)\lambda$. It never reaches $\mu$. A refinement loop using only the number of terms can therefore run forever, even though every finite step is mathematically correct. In the full surreal fine topology, even stronger nonconvergence phenomena occur; the manuscript explicitly distinguishes it from a fixed workspace's valuation topology \cite{SourceAnalysis}.

\subsection{Rounding and integer-part conventions}
For an infinite surreal there is no ordinary integer floor in $\Z$. A package must state what its target ``integers'' are before defining a global floor. One useful normal-form choice is the ordered subring $\Pi\oplus\Z$, where $\Pi$ is the purely infinite additive part. Products of purely infinite parts have only negative exponents, so this is indeed a subring.

For $x=p+r+\eta$, with $p$ purely infinite, $r\in\R$, and $\eta$ infinitesimal, define
\[
 \lfloor x\rfloor_{\Pi\oplus\Z}
 =p+\begin{cases}
 \lfloor r\rfloor-1,&r\in\Z\text{ and }\eta<0,\\
 \lfloor r\rfloor,&\text{otherwise}.
 \end{cases}
\]
It is the unique element $m\in\Pi\oplus\Z$ such that $m\le x<m+1$: noninteger $r$ has positive ordinary distance from its neighboring integers, while at an integer the infinitesimal sign selects the side. Any distinct purely infinite part differs by an infinite quantity and cannot satisfy the same unit interval condition.

Thus $\lfloor1-t\rfloor=0$, whereas $\lfloor\omega/2\rfloor_{\Pi\oplus\Z}=\omega/2$. The second answer is not an ordinary integer. This operation is effective only when the backend can split the normal form, decide the required ordinary integer predicate, and determine the infinitesimal sign. A compact rational node need not expose its entire purely infinite part as a finite list.

\subsection{A limit is not substitution of an infinite number}
An ordinary real limit and a surreal evaluation should be different commands. For example,
\[
 \lim_{X\to+\infty}\frac1{X+1}=0\quad\text{in ordinary real analysis},
 \qquad \frac1{\omega+1}>0\quad\text{in }\No.
\]
The latter retains the infinitesimal information discarded by the former. Similarly, the ordinary sequence $1/n$ does not converge finely to zero in the full surreal field: a fixed positive surreal infinitesimal is smaller than every term. Limit modes, Hahn summation, and finite-scale specialization must therefore remain explicit operations, rather than different spellings of one implicit ``infinite evaluation'' routine \cite{SourceAnalysis}.

\subsection{Safe numerical views}
There is no ordinary floating-point value of $\omega$. A useful numerical interface instead offers coefficient approximation, standard-part approximation for certified finite inputs, or an explicit finite-scale specialization. The specialization $t_1=10^{-6}$ can illustrate a one-scale expression, but it is not its surreal value.

For several scales, numerical choices can respect any prescribed finite collection of hierarchy inequalities, not the entire statement $t_2<t_1^n$ for all ordinary $n$. Results from such a specialization must be labeled as simulations. They cannot prove exact non-Archimedean order or resolve an expression whose leading coefficient is symbolically unknown.

\section{Polynomial algebra, root lifting, and finite geometry}\label{sec:algebra}
\subsection{Finite polynomial problems}
Once a coefficient backend provides exact field operations and zero tests, ordinary finite-degree polynomial arithmetic, Euclidean division, greatest common divisors, resultants, and finite-dimensional linear algebra can be implemented over it. Factorization and root finding need the additional capabilities of the particular field; they are not consequences of having a multiplication method alone.

The real-closed core $R_d$ and its complexification $C_d$ from \cref{sec:core} give an especially strong setting. Gaussian elimination decides matrix rank and solves consistent systems. Characteristic polynomials and algebraic root selectors represent eigenvalues. Positive square roots give distances, and exact determinant signs decide orientation. The trigonometry manuscript's triangle identities can thus be supported without infinite circular phases: normalized directions lie on the finite-coordinate unit circle \cite{SourceTrig}.

For coefficients represented by arbitrary analytic nodes, these algorithms become conditional on zero testing. A determinant that fails to simplify is not a proof of nonsingularity. A domain adapter should expose the conditions under which a matrix result was obtained.

\subsection{Valuations of roots and a Newton step}
For $P(X)=\sum_{j=0}^m a_jX^j$, suppose a nonzero root begins with $c\,t^\lambda$. The candidate leading valuation must satisfy
\begin{equation}\label{eq:newtonmin}
 \min_j\bigl(v(a_j)+j\lambda\bigr)
 \quad\text{is attained at least twice}.
\end{equation}
Otherwise the unique term of least valuation could not cancel. After selecting a candidate $\lambda$, the coefficients attaining the minimum give the initial equation
\[
 \sum_{j\in J}\lc(a_j)c^j=0.
\]
Its nonzero roots provide candidate leading coefficients. Substituting $X=t^\lambda(c+Y)$ and normalizing produces a smaller perturbation problem. Fractional $\lambda$ can force ramification. These are the algebraic steps behind Newton--Puiseux and generalized Newton-polygon methods; grid-based treatments supply a developed framework \cite{vdH}.

The following example needs no infinite search:
\[
 X^2-2X+1-t=0
 \quad\Longrightarrow\quad X=1\pm t^{1/2}.
\]
Both roots have standard part $1$, but they are distinct exact elements. Storing only their standard parts loses multiplicity and branch information. Similarly $X^2+t=0$ has roots $\pm i t^{1/2}$, not real roots with uncertain sign.

For a simple normalized root, one can also solve coefficients recursively or use Newton corrections
\[
 x_{k+1}=x_k-\frac{P(x_k)}{P'(x_k)}.
\]
The actual implementation must justify invertibility of $P'(x_k)$ and the precision improvement. In a higher-rank field, repeatedly doubling a positive error valuation need not reach a cutoff in a strictly higher Archimedean class. Newton iteration is therefore a coefficient-refinement method with a stated target policy, not a universal convergence proof in the full surreal topology.

\subsection{Implicit analytic roots}
Suppose an ordinary analytic equation has $f(c)=0$ and $a=f'(c)\ne0$, and consider
\[
 f\sh(z)+\tau g\sh(z)=0
\]
for a positive infinitesimal $\tau$. The supplied analysis manuscript gives the unique solution in the monad of $c$ as a support-controlled implicit germ. Its first coefficients are
\begin{equation}\label{eq:implicit}
 z=c-\frac{g(c)}a\tau+
 \left(\frac{g(c)g'(c)}{a^2}
       -\frac{f''(c)g(c)^2}{2a^3}\right)\tau^2+O(\tau^3).
\end{equation}
A CAS can store the implicit equation, selected monad, and nonzero-Jacobian certificate as the exact answer, then compute coefficients on demand. The proof is coefficient recursion, not a numerical root search in a large finite approximation \cite[``Computing a lifted simple root'']{SourceAnalysis}.

At multiple roots, the correct interface may return a finite cluster object rather than arbitrarily numbered individual branches. The manuscripts' preparation and coprime-factor lifting results retain the support of \emph{factor coefficients}, while individual roots can require fractional scales. An implementation can exploit that distinction: factor coefficients are often stable under perturbation even when individual root labels are not.

\subsection{Near-tangency conditioning is algebraic information}
The trigonometry manuscript proves, for $0<\delta\in\mm$,
\begin{equation}\label{eq:acos}
 q(\delta)=\arccos(1-\delta)
 =\sqrt{2\delta}\left(1+\frac\delta{12}
       +\frac{3\delta^2}{160}+\frac{5\delta^3}{896}+O(\delta^4)\right).
\end{equation}
If $e\ne0$ and $e\prec\delta$, its branch stability formula gives
\begin{equation}\label{eq:acosprecision}
 v\bigl(q(\delta+e)-q(\delta)\bigr)
       =v(e)-\tfrac12v(\delta).
\end{equation}
The inverse derivative therefore loses half the valuation of the tangency parameter. This is not a roundoff anecdote: it is an exact input--output precision rule that a branch-aware algebraic or trigonometric evaluator can propagate \cite[``Tangency, angular splitting, and loss of valuation'']{SourceTrig}.

\section{Ordinary germs and exact infinitesimal special-function values}\label{sec:germs}
\subsection{The canonical local lift}
For an ordinary holomorphic germ $f$ at $c\in\C$, and $\eta\in\mm$, define
\begin{equation}\label{eq:taylorlift}
 f\sh(c+\eta)=\sum_{n\ge0}\frac{f^{(n)}(c)}{n!}\eta^n.
\end{equation}
The union of supports lies in the positive monoid generated by $\supp(\eta)$, so the sum is strongly summable. Ordinary values such as $f(c)$ are evaluated or named first. It is not correct to insert a nonzero ordinary $c$ directly into a Hahn power series about zero and invoke ordinary convergence.

\begin{proposition}[A reusable exact analytic constructor]\label{prop:lift}
Suppose a function descriptor specifies an ordinary analytic germ, an exact center, and an effective Taylor-coefficient procedure. Given an exact infinitesimal argument with a certified support representation, \eqref{eq:taylorlift} supplies an exact symbolic value. If the support and coefficient dependencies are effective, each requested finite coefficient calculation terminates. Equality of arbitrary resulting values is not guaranteed by this construction alone.
\end{proposition}
\begin{proof}
The positive-monoid support lemma supplies semantic existence. The descriptor fixes every coefficient, hence a unique value. Each specified effective coefficient request uses finitely many certified dependencies. The final qualification follows from \cref{thm:halting}: a general formal-series constructor cannot automatically acquire a complete zero test.
\end{proof}
The source manuscripts prove that the lift respects analytic identities, permitted composition, and differentiation. Thus the CAS can reuse an ordinary identity after checking that its germs and branch choices agree on a common domain \cite{SourceAnalysis,SourceTrig}.

\subsection{Which familiar functions can be included?}
At regular ordinary centers, this mechanism applies to elementary analytic functions and to specified branches of special functions: exponential, logarithm away from its selected cut, trigonometric functions, gamma away from its poles, error functions, Airy functions, Bessel functions with their parameter and branch conditions, and hypergeometric functions at regular parameter and argument values. This is a rule about known ordinary germs, not an assertion that every possible parameter value or singularity is regular.

Ordinary CAS differentiation and series tools are valuable coefficient engines. Wolfram Language supplies exact special-function expressions, derivatives, and local expansions \cite{WLSpecial}. A successful coefficient call does not, by itself, certify an entire surreal support or a global function domain. The extension layer must preserve that missing information.

For example, $\Gamma\sh(2+t)$ is an exact value given by the ordinary germ at $2$. Coefficients may involve constants such as Euler's constant and derivatives of gamma. They lie outside the real-algebraic coefficient core. The system should promote to an exact ordinary-constant domain with correspondingly weaker general equality guarantees, not silently call those constants algebraic.

\subsection{Singular centers require a different constructor}
At a pole one uses a finite-principal-part Laurent germ, on a punctured infinitesimal neighborhood. At an algebraic branch point one uses a ramified coordinate and a selected Puiseux branch. Formula \eqref{eq:acos} is a basic example: ordinary Taylor expansion of $\arccos$ at $1$ is not an analytic power series in $\delta$, but a convergent ordinary series after the positive square-root uniformizer is introduced.

At an essential singularity, a Laurent expression can fail Hahn summability. If $z$ is nonzero infinitesimal, the terms $z^{-n}/n!$ have decreasing valuations, so the ordinary series for $e^{1/z}$ cannot be used as a Hahn sum. The source analysis manuscript explicitly distinguishes this failed substitution from a separately selected global exponential evaluated at $1/z$ \cite[``Essential singularities and a concrete summability boundary'']{SourceAnalysis}.

\subsection{Formal, analytic, and asymptotic descriptors are different}
A formal series with computable coefficients can be meaningful at infinitesimal inputs even when it is not an ordinary analytic germ. An asymptotic expansion of an ordinary function is yet another kind of data. In particular, a power-scale expansion at $+\infty$ cannot distinguish an ordinary function $f(X)$ from $f(X)+e^{-X}$ when the added term is beyond every inverse-power order.

Under Gonshor semantics, the corresponding difference at $\omega$ is the positive surreal $\ExpNo(-\omega)$, not zero. An exact extension of a special function at infinity must therefore state how exponentially small data, sectors, branches, and any summation choices are fixed. The constructive extension and integration theory of Costin--Ehrlich is important established background for a large class of resurgent functions; it is not identical to the elementary local Taylor-lift interface or to the manuscript's contour functional \cite{CostinEhrlich}.

\section{Real exponentials, logarithms, and transserial growth}\label{sec:exp}
\subsection{Canonical real semantics, partial computational backends}
The established surreal exponential $\ExpNo:\No\to\No_{>0}$ and its inverse $\LogNo$ are canonical choices in the supplied framework. They extend ordinary real exponential and logarithm and satisfy the ordered group laws \cite{Gonshor,vdDE}. Thus a finite node $\code{ExpNo[x]}$ can have an exact meaning whenever $x$ does, even when a particular series backend cannot normalize its value.

Writing $x=p+r+\eta$, with $p$ purely infinite, $r\in\R$, and $\eta$ infinitesimal, gives
\begin{equation}\label{eq:expsplit}
 \ExpNo(x)=\ExpNo(p)e^r\Expi(\eta),\qquad
 \Expi(\eta)=\sum_{n\ge0}\frac{\eta^n}{n!}.
\end{equation}
The infinitesimal factor is directly computable by support-controlled expansion. The purely infinite factor requires Gonshor-compatible monomial/transseries machinery or must remain an exact semantic node. Defining it by $\sum p^n/n!$ would be invalid: for nonzero purely infinite $p$, term valuations decrease.

For a positive normal form $a=c\,t^\gamma(1+\eta)$, $c>0$, one similarly has
\begin{equation}\label{eq:logsplit}
 \LogNo(a)=\log c+\LogNo(t^\gamma)+\Logi(1+\eta),
 \qquad \Logi(1+\eta)=\sum_{n\ge1}\frac{(-1)^{n+1}\eta^n}{n}.
\end{equation}
The middle term is \emph{not} a license to conflate the Conway monomial map with arbitrary exponential powers. A backend must implement the applicable monomial-logarithm identities or retain that term as an atom.

\subsection{An exact exp--log expression language}
One practical language starts with a decidable algebraic core and permits $\ExpNo$, $\LogNo$ on certified positive arguments, and positive-base real powers. Its expression graphs are finite and denotationally exact. It can simplify group identities, cancel inverse functions on their domains, extract known leading scales, and perform infinitesimal expansions.

It should not claim that every equality between arbitrary such graphs has a known terminating decision procedure. In particular, broadening the ordinary constant domain already complicates equality before any surreal exponent is introduced. A partial simplifier can return a smaller equivalent graph and a proof of its rewrites without producing a canonical normal form for the entire language.

The backend should distinguish a \emph{structural constructor} from a \emph{normalization request}. The former can often succeed immediately. The latter may return a normalized object, a conditional result, a residual exact node, or a resource-limit result. Conflating these operations either rejects useful exact expressions unnecessarily or encourages unsupported canonicalization claims.

\subsection{Implicit functions can add exact infinite values}
There are useful exact constructors beyond a fixed list of elementary functions. For example, a positive-real Lambert function can be specified by
\[
 W_+(y)\ge0,\qquad W_+(y)\ExpNo(W_+(y))=y,\qquad y\ge0.
\]
The ordinary real existence statement is first order in the ordered exponential language, so the elementary-extension result supplies existence for nonnegative surreal $y$ as well \cite{vdDE}. Uniqueness follows directly: the map $x\mapsto x\ExpNo(x)$ is strictly increasing for $x\ge0$, by positivity and the strict increase of the real exponential.

Consequently a selected implicit-root node gives an exact value such as $W_+(\omega)$, even before a normalizer has computed its logarithmic-scale expansion. A transseries inversion algorithm can provide that expansion within its supported class. The existence and uniqueness argument is not a generic transcendental root-finding algorithm: it is a certificate for this particular function. General equations without comparable certificates remain conditional or partial.

\subsection{A transseries backend is the natural larger project}
An exponential--logarithmic monomial hierarchy, small-monomial grids, coefficient algorithms, and support-preserving substitutions provide a more systematic backend than ad hoc rewrite rules. It can expose formal differentiation, integration within its supported differential class, composition for admissible large positive arguments, and functional inversion under the appropriate hypotheses. Such operations have extensive mathematical antecedents in the transseries literature \cite{vdH}.

The embedding into surreal numbers matters: it identifies the formal large variable with $\omega$ and specifies the exponential structure. A set of independent formal labels called $\log\omega$ and $e^\omega$ without their structural relations is not the same ordered exponential field.

Recent work continues to strengthen this framework. Mantova's 2026 preprint, in its May 9 revision, establishes monotonicity of omega-series composition in its second argument and a Taylor approximation theorem, with an intermediate-value consequence for the interpreted functions \cite{Mantova2026}. These are semantic theorems for specified transserial classes, not a claim that arbitrary CAS expressions or arbitrary coefficient programs can be normalized. Likewise, hyperseries are a longer-term representational direction rather than a prerequisite for a first useful implementation \cite{Hyperseries}.

\section{Surcomplex arithmetic, angles, logarithms, and powers}\label{sec:complex}
\subsection{Pairs give a robust exact representation}
Represent $z=x+iy$ as a pair of real-surreal objects with a common compatible domain. Then
\begin{align*}
 (x+iy)(u+iv)&=(xu-yv)+i(xv+yu),\\
 \overline z&=x-iy,\qquad |z|^2=x^2+y^2,\\
 z^{-1}&=\frac{x-iy}{x^2+y^2}\qquad(z\ne0).
\end{align*}
The squared norm lies in the real base field; the norm $|z|=\sqrt{x^2+y^2}$ may require its real closure. No ordering should be imposed on the complexification as a field. Comparisons must specify real part, modulus, valuation, or another intended relation.

\subsection{All directions admit finite surreal angles}
The trigonometry manuscript proves that the unit circle
\[
 \mathbb T(\No)=\{u\in\SC:u\overline u=1\}
\]
is the image of canonical $\cis$ on finite real surreals, with kernel $2\pi\Z$. Thus every nonzero surcomplex number has a modulus and a finite angle \emph{class}, even when its coordinates or modulus are infinite \cite[``Polar decomposition and the angle group'']{SourceTrig}.

For computation, let $u=z/|z|$ and $u_0=\st(u)$. Then
\begin{equation}\label{eq:anglesplit}
 u=u_0\Expi(i\delta),\qquad
 \delta=-i\Logi(u/u_0)\in\mm\cap\No.
\end{equation}
The pair $(u_0,\delta)$ avoids an arbitrary angle cut. It is an exact phase representation, and it separates the ordinary unit direction from the infinitesimal correction. Extracting $u_0$ still requires the scalar backend's standard-part capability.

An entirely algebraic alternative, avoiding angle evaluation, is the Cayley coordinate
\begin{equation}\label{eq:cayley}
 u(q)=\frac{1+iq}{1-iq}
      =\frac{1-q^2}{1+q^2}+i\frac{2q}{1+q^2},\qquad q\in\No.
\end{equation}
It parametrizes the circle except $-1$; that omitted point must be checked separately in solving and integration. Infinite $q$ is allowed and describes an infinitesimal neighborhood of the missing direction.

\subsection{Canonical domains for exponential and trigonometry}
Write
\[
 \mathscr H=\{x+iy:y\text{ finite},\ x\in\No\},\qquad
 \mathscr V=\{x+iy:x\text{ finite},\ y\in\No\}.
\]
The manuscript's canonical strip exponential is
\begin{equation}\label{eq:stripexp}
 E_{\mathscr H}(x+iy)=\ExpNo(x)\cis(y).
\end{equation}
It is onto $\SC^\times$ and has kernel $2\pi i\Z$. On $\mathscr V$, canonical sine and cosine are
\begin{align}
 \Sin(x+iy)&=\sin x\cosh y+i\cos x\sinh y,\\
 \Cos(x+iy)&=\cos x\cosh y-i\sin x\sinh y.
 \label{eq:striptrig}
\end{align}
Here the real hyperbolic functions are defined by $\ExpNo(\pm y)$. These formulas admit exact symbolic nodes even at infinite $y$, conditional only on the scalar backend capabilities. In particular,
\[
 \Sin(i\omega)=i\sinh\omega
\]
is canonical in this framework. By contrast, $\sin\omega$ and $E(i\omega)$ are not fixed by its canonical finite-angle data. The distinction is about the real circular phase, not merely about whether the argument is large \cite[\S10]{SourceTrig}.

For finite real $r+\eta$, evaluate the ordinary constants first:
\begin{align}
 \sin(r+\eta)&=\sin r\sum_{n\ge0}\frac{(-1)^n\eta^{2n}}{(2n)!}
       +\cos r\sum_{n\ge0}\frac{(-1)^n\eta^{2n+1}}{(2n+1)!},\\
 \cos(r+\eta)&=\cos r\sum_{n\ge0}\frac{(-1)^n\eta^{2n}}{(2n)!}
       -\sin r\sum_{n\ge0}\frac{(-1)^n\eta^{2n+1}}{(2n+1)!}.
\end{align}
The inverse functions have the natural real-surreal domains $\arctan:\No\to(-\pi/2,\pi/2)$, $\arcsin:[-1,1]\to[-\pi/2,\pi/2]$, and $\arccos:[-1,1]\to[0,\pi]$. For positive infinite $x$, use
\[
 \arctan x=\frac\pi2-\arctan(1/x),
\]
so an exact finite-angle value is available without defining circular sine at $x$ itself.

\subsection{Two logarithm branches must not share one cache key}
An actual principal-angle convention chooses a representative in $(-\pi,\pi]$. The analysis manuscript instead also defines a \emph{standard-direction-anchored} logarithm
\begin{equation}\label{eq:shadowlog}
 \mathcal L(z)=\LogNo|z|+i\Arg(u_0)+\Logi(u/u_0),
 \qquad u=z/|z|,\quad u_0=\st(u),
\end{equation}
where $\Arg(u_0)$ is an ordinary principal argument. This is fine-locally analytic on the full punctured class plane, with finite imaginary part, but is not a coherent primitive of $1/z$ around the ordinary puncture \cite[``A global fine-analytic logarithm'']{SourceAnalysis}.

These two choices differ at infinitesimal crossings of the ordinary branch cut. For $\eta>0$ infinitesimal and $z=-1-i\eta$, the actual principal angle is
\[
 -\pi+\arctan\eta,
\]
whereas the anchored angle is $\pi+\arctan\eta$. Their difference is $2\pi$. The latter is permitted to lie infinitesimally above $\pi$. It must not be silently reduced back into the principal interval while retaining a claim of the same global fine-local analytic behavior.

A function descriptor therefore needs a branch policy such as \code{ActualPrincipal} or \code{StandardDirectionAnchored}. The chosen policy belongs to exact expression identity, domain checks, differentiation, and serialization.

\subsection{General complex powers have additional domain obligations}
For a selected logarithm $L$, the expression
\[
 z^w=E_{\mathscr H}(wL(z))
\]
is canonical by this formula only when $wL(z)\in\mathscr H$. Its imaginary part need not be finite for arbitrary surcomplex $w$. Outside that strip one must select a global exponential extension or retain an explicitly unresolved power node.

Integer powers need none of this and should remain algebraic. Ordinary rational powers can also be implemented by algebraic root selection and finite-angle branch rules rather than unnecessarily introducing an unrestricted exponential. For positive real base and real surreal exponent, the real definition $\ExpNo(w\LogNo z)$ is available. A single unqualified \code{Power} rewrite table cannot safely cover all these cases.

\section{Infinite phases and exact trigonometric solving}\label{sec:phases}
\subsection{A global extension is a named mathematical structure}
Let $\Pi$ be the additive class of purely infinite real surreals. Normal forms give the additive splitting
\[
 \No=\Pi\oplus\OO_{\R},\qquad x=p+\fin(x).
\]
The trigonometry manuscript classifies extensions of the circular group law by characters
\[
 \chi:(\Pi,+)\longrightarrow(\mathbb T(\No),\cdot).
\]
For a specified character, define
\begin{equation}\label{eq:character}
 U_\chi(p+u)=\chi(p)\cis u,\qquad u\text{ finite}.
\end{equation}
Its real and imaginary parts give global cosine and sine satisfying the addition laws and the manuscript's fine-local differential identities. Conversely every extension with this group-law property has that form \cite[``All circular group-law extensions'']{SourceTrig}.

The convention $\chi=1$ gives $U_0(x)=\cis(\fin(x))$, hence $\sin_0\omega=0$. It is an available convention, not an identity forced by the finite-angle theory. Other characters give different phases at $\omega$. An exact CAS should default to \code{ExtensionRequired} for such a circular phase unless the user selects a provider. It should not silently adopt the trivial character.

A provider must respect additivity across its domain. Independently assigning phases to arbitrary related inputs can violate the group law. For instance, the phase at $2p$ must be the square of the phase at $p$, and a phase assigned at $p/2$ must be a compatible square root. A character descriptor is therefore structural data, not a mutable table of unrelated evaluations.

\begin{proposition}[Why global extensions have new periods]\label{prop:periods}
Every extension \eqref{eq:character} has an infinite period in each nonzero purely infinite coset $p+\OO_\R$.
\end{proposition}
\begin{proof}
Finite-angle $\cis$ is already onto the entire surreal unit circle. Choose finite $a$ with $\cis a=\chi(p)^{-1}$. Then $U_\chi(p+a)=1$. If $p\ne0$, the element $p+a$ is infinite, and the homomorphism law makes it a period. This reproduces the obstruction established in the supplied trigonometry manuscript.
\end{proof}
Thus a global phase provider must not also advertise that its only periods are ordinary integer multiples of $2\pi$. Those two requirements are incompatible in this framework. Exact results computed with different providers must retain their provider identifiers, even if they happen to agree on all ordinary test values.

\subsection{Finite trigonometric polynomials reduce to algebra}
Let $n$ be an ordinary finite integer and consider, on $\mathscr V$,
\[
 T(z)=\sum_{k=-n}^n c_k E_{\mathscr H}(ikz),\qquad c_k\in\SC.
\]
Setting $U=E_{\mathscr H}(iz)$ gives
\begin{equation}\label{eq:trigalgebra}
 T(z)=U^{-n}P(U),\qquad
 P(U)=\sum_{k=-n}^n c_kU^{k+n}.
\end{equation}
The manuscript proves that nonzero roots of $P$ correspond, with multiplicity, to zeros of $T$ modulo $2\pi$ in the canonical strip. The number is at most $2n$, and is exactly $2n$ when both extreme coefficients are nonzero \cite[``Algebraization and the $2n$ root bound'']{SourceTrig}.

This gives an implementable exact solver. Build $P$ with the scalar polynomial backend; solve it in the appropriate complex algebraic extension; discard $U=0$; retain multiplicities; convert a nonzero $U$ to a strip angle class through a logarithm. This final analytic reconstruction may leave the algebraic coefficient core, even when solving for $U$ stays inside it. To select finite real angles, impose $U\overline U=1$. Alternatively, use the Cayley parametrization and a real algebraic solver, remembering the omitted point $-1$.

An ``infinite integer frequency'' is not an ordinary polynomial degree. Replacing $n$ by $\omega$ in this algorithm is meaningless. Infinite-frequency oscillation requires a different function class and, in general, an infinite-phase provider.

\subsection{Positivity must see infinitesimal angles}
For positive infinitesimal $\varepsilon$, the source's example
\[
 T(\theta)=(\cos\theta-\varepsilon)^2-\varepsilon^4
\]
is positive at every ordinary real angle but is negative at the finite surreal angle $\theta=\arccos\varepsilon$. A positivity test using only ordinary angular samples or standard parts is therefore unsound \cite[``Ordinary-angle testing misses infinitesimal negativity'']{SourceTrig}.

For coefficients in the real-closed algebraic core, one can instead reduce circle positivity to a semialgebraic question using \eqref{eq:cayley}, and separately inspect the omitted circle point. The source's Fej\'er--Riesz theorem also provides a finite algebraic factorization route. Both approaches test the actual ordered field, not only its ordinary shadow.

\section{Differentiation, composition, integration, and function types}\label{sec:calculus}
\subsection{Three derivatives with different answers}
A system needs to distinguish differentiation of a function in a variable from differentiation of its scalar coefficients as transseries. For a function $F(z)$ with surreal parameters, the source's derivative $D_z$ holds every scalar parameter fixed. Thus $D_z\omega=0$ and
\[
 D_z(\omega z^2+t\sin z)=2\omega z+t\cos z.
\]
A formal series derivative $D_t$ instead satisfies $D_tt=1$. A selected surreal-field derivation $\partial$ may satisfy $\partial\omega=1$; the normalized Berarducci--Mantova structure is an important established example \cite{BMDerivations}. Then, for $t=\omega^{-1}$,
\[
 \partial t=-t^2,\qquad D_t(t^3)=3t^2,
 \qquad \partial(t^3)=-3t^4.
\]
These are not contradictory outputs: they are different operators.

In particular, $D_z$ commutes with coefficientwise Hahn assembly while holding monomials constant, as in the analysis manuscript. Extending a scalar derivation through a Hahn sum requires its own strong-additivity hypotheses. For general monomials with nonconstant surreal exponents, one must use the chosen derivation's rules, not an unqualified power rule copied from ordinary integer exponents.

A reasonable API therefore uses \code{FunctionDerivative}, \code{FormalSeriesDerivative}, and \code{ScalarDerivation}, with an explicit derivation identifier in the last case. Converting between these interpretations should be a deliberate operation.

\subsection{Composition has several distinct contracts}
Formal power-series composition $A(B(X))$ is immediately defined when the inner constant term is zero, because each resulting coefficient depends on finitely many terms. A nonzero inner constant generally requires an analytic re-expansion or a separate summability condition. At finite surreal inputs, the ordinary-center Taylor lift supplies that re-expansion when the ordinary analytic domains match.

For a Hahn-coherent datum, the analysis manuscript allows composition through an inner map whose ordinary leading coefficient lands in the outer ordinary domain and whose remaining terms have positive support. This common-domain condition is stronger than pointwise existence. For a transseries function at infinity, admissible composition uses its large-argument theory instead. These contracts should not be collapsed into a single unchecked substitution routine \cite{SourceAnalysis,BMFunctions}.

\subsection{Local primitives and finite meromorphic residues}
A formal analytic germ $\sum a_nX^n$ has the normalized formal primitive
\[
 \sum_{n\ge0}\frac{a_n}{n+1}X^{n+1}.
\]
It can be represented by a new coefficient generator with the same local semantic obligations. A Laurent germ with finite principal part is handled termwise, except that an $X^{-1}$ term requires a logarithm descriptor or remains a formal logarithmic primitive. The residue is simply the coefficient of $X^{-1}$.

For rational functions over $C_d$, polynomial division and finite partial fractions give effective residue computations once algebraic roots and multiplicities are available. An antiderivative need not remain rational; logarithms and their branch information must be added. In a broader analytic coefficient domain, a definite integral may be stored exactly without a closed-form reduction. Existence, representability by an integral node, and reduction to elementary functions are separate outcomes.

\subsection{Hahn-coherent functions and contour functionals}
Following the supplied analysis manuscript, a coherent datum on an ordinary domain $U\subseteq\C$ has the form
\begin{equation}\label{eq:coherent}
 F=\sum_{\gamma\in S}f_\gamma t^\gamma,
 \qquad f_\gamma\in\OO(U),
\end{equation}
where $\OO(U)$ denotes the ordinary holomorphic functions on $U$, with one common well-ordered support. Its evaluation at $c+\eta\in U\sh$ is
\[
 F\sh(c+\eta)=\sum_{\gamma,n}
       \frac{f_\gamma^{(n)}(c)}{n!}t^\gamma\eta^n.
\]
A computational object should therefore include both a support description and an ordinary analytic domain descriptor. Merely holding an unrelated germ at each center does not provide the coherent data required by the source's identity and contour theorems.

For an ordinary piecewise smooth path $C$ in $U$, the manuscript defines
\begin{equation}\label{eq:contour}
 \HInt_C F(z)\,dz
   =\sum_\gamma t^\gamma\int_C f_\gamma(z)\,dz.
\end{equation}
The resulting support is contained in $S$. Each ordinary coefficient integral can be evaluated by the underlying CAS when possible, or retained as an exact integral descriptor. An effective infinite coefficient family again needs an effective coefficient procedure. Nonstandard endpoints add the source's Taylor corrections to ordinary endpoint integrals.

This is a viable exact symbolic integration layer. It is not justified by taking an ordinary sequence of Riemann sums in the full surreal fine topology. That topology does not make ordinary nonconstant paths into fine-continuous paths; the source explicitly separates the contour functional from such an interpretation \cite[\S6]{SourceAnalysis}.

\subsection{A complete cancellation of infinite residues}
The manuscript's example
\[
 M(z)=\frac{e^z}{z^2-t},\qquad s=t^{1/2},
\]
has poles $\pm s$ and residues $e^s/(2s)$ and $-e^{-s}/(2s)$. Each residue is infinite. Their sum is nevertheless the exact finite value
\begin{equation}\label{eq:residue}
 \frac1{2\pi i}\HInt_{|z|=1}M(z)\,dz
  =\frac{\sinh s}{s}
  =\sum_{n\ge0}\frac{t^n}{(2n+1)!}
  =1+\frac t6+\frac{t^2}{120}+\cdots.
\end{equation}
Here the exponential on the infinitesimal poles is the canonical lift, and the contour is the source's ordinary-circle coefficientwise functional \cite[``An explicit cancellation of infinite residues'']{SourceAnalysis}.

A correct CAS should combine the exact residue expressions before applying standard part or coarse truncation. Treating the individual residues as generic infinities would produce an indeterminate expression and destroy a computable answer. Keeping their algebraic dependencies makes the cancellation automatic.

\subsection{All-scale coherence is too strong for a general special-function type}
The source analysis manuscript proves that functions coherent in its prescribed sense at every surreal scale are precisely polynomials, and that its all-scale meromorphic fractions are precisely rational functions. These statements use the full class of scales and the manuscript's uniform-support hypothesis \cite[``Affine scales and global rigidity'']{SourceAnalysis}.

A global special-function object should therefore not automatically claim all-scale Hahn coherence. It can instead advertise local germs, selected coherent charts, a transserial interpretation at infinity, or a sectorial extension provider. The compatible charts and transition rules are part of the object. This preserves the source's positive theorems without imposing a category that excludes the intended transcendental functions.

\section{Wolfram Language integration without semantic overloading}\label{sec:wl}
\subsection{Reuse the host CAS as a coefficient and finite-algebra engine}
The strongest immediate reuse is exact ordinary arithmetic, polynomial normalization, algebraic-number arithmetic, differentiation of coefficient functions, and local series generation. These are mature host operations. A surreal extension should adapt their inputs and outputs rather than ask the host to infer non-Archimedean semantics from ordinary symbols.

For example, use ordinary \code{Together} and \code{Cancel} on polynomial fractions after their variables have been assigned a separate monomial-domain interpretation. Use ordinary exact algebraic numbers for the coefficient layer. Perform geometric operations over the custom real-closed extension, not by substituting \code{Infinity} into real or complex numeric routines.

\subsection{What documented built-in representations do not establish}
The documented \code{SeriesData} representation stores a finite coefficient list on a rational-power lattice together with an omitted-order term. It is well suited to a one-parameter finite-jet adapter. Its \code{Normal} conversion removes the order term; it does not prove that the remainder is zero. It is not, by its documented contract, an arbitrary well-ordered Hahn support or an exact infinite-series generator \cite{WLSeries}.

The current \code{Root} documentation includes both indexed polynomial-root representations and neighborhood representations for general equations, including transcendental equations. It would therefore be inaccurate to describe every current \code{Root} object as algebraic over $\Q$. Nevertheless, those documented representations do not themselves specify root isolation in arbitrary surreal-valued coefficient fields. A custom algebraic surreal root descriptor still needs its ordered-field selector and domain information \cite{WLRoot}.

Ordinary asymptotic tools can generate valuable candidate expansions, but an asymptotic expansion is not automatically an exact surreal extension. As a comparison, Sage's asymptotic-ring framework also explicitly works with growth terms and error terms, rather than purporting to encode every surreal \cite{SageAsymptotic}. These systems are sources of algorithms and interfaces, not substitutes for the semantic choices made here.

\subsection{Proposed object layers}
The following names are \emph{proposed interfaces}, not claims about existing Wolfram built-ins:
\begin{lstlisting}
SurrealDomain[coefficientDomain, exponentDomain, interpretation]
SurrealRational[domain, numerator, denominator]
SurrealAlgebraic[domain, polynomial, rootSelector]
CertifiedHahn[domain, supportDescriptor, coefficientDescriptor]
SurrealExpression[domain, expressionGraph]
SurrealJet[domain, center, remainderDescriptor]
Surcomplex[realPart, imaginaryPart]
AnalyticGerm[functionDescriptor, center, branch, domain]
CoherentFunction[ordinaryDomain, support, coefficients]
\end{lstlisting}
In actual code, a domain identifier can point to an immutable association or a validated expression tree. Constructors check invariants before producing objects. Promoting a rational node to a complex pair or an algebraic extension must preserve the embedding, not simply discard the older domain tag.

Each operation reports the capability actually used. A normalization can succeed in the rational backend while leaving one transcendental coefficient symbolic. A positivity query can succeed from a certified leading coefficient without requiring full normalization. A valuation cutoff can produce a compressed exact block rather than a flat list. These are useful, distinct successful results.

\subsection{Evaluation control and conservative dispatch}
Inert symbolic heads and carefully controlled evaluation keep the host from simplifying away a domain distinction before the extension sees it. Wolfram's \code{Inactive} and upvalue mechanisms offer tools for such an integration \cite{WLInactive,WLUpValues}. However, the first reference implementation should use explicit operations such as \code{SRAdd} and \code{SRSin}. It is easier to test their domains than to debug an early collection of global automatic rewrites.

Once the core is stable, narrowly targeted upvalues can provide natural notation for operations whose semantics are unconditional in the operands' domain. Do not globally unprotect \code{Plus}, \code{Power}, \code{Sin}, or \code{Log} and redefine them for arbitrary expressions. Branch-dependent transformations should remain explicit or return their conditions.

Likewise, \code{NumericQ} is not a declaration that an object belongs to a non-Archimedean ordered field. Its documented role concerns numeric expressions \cite{WLNumericQ}. Marking an infinite surreal as an ordinary numeric value can cause unrelated numerical algorithms to attempt impossible real approximation. Use a custom exact-domain predicate; supply numerical views only through explicitly described coefficient approximation or finite-scale specialization.

\subsection{Results, failures, and unresolved obligations}
A proposed result protocol distinguishes \code{Exact}, \code{CertifiedApproximation}, \code{Conditional}, \code{Unknown}, and \code{Failure}. A missing zero proof is not a mathematical domain error. An input outside the chosen sine strip is not a proof that no extension exists. Division by a certified zero is an error, whereas division by an object of unknown zero status creates a nonzero obligation.

Examples of useful obligations are: positive argument for real logarithm; finite real part for canonical strip sine; positive support for an infinitesimal series substitution; a root selector valid in the current ordered field; common ordinary domain for coherent composition; and a compatible phase character for infinite circular input. An implementation can attach these obligations to an exact conditional node rather than prematurely refusing all further algebra.

\subsection{Proof-carrying expressions without pretending to have a proof assistant}
A practical certificate can be a small derivation tree using trusted constructors: finite support; positive-support monoid; exact polynomial identity; a host algebraic-number sign certificate; or a branch identity under recorded inequalities. The checker should be smaller and more deterministic than the expression-search engine.

This does not amount to a machine-checked formalization of the entire surreal class. If certificates rely on trusted mathematical theorems from the manuscripts, that trust boundary should be documented. A later formal checker could verify the finite algebraic kernel and the soundness of selected support constructors without requiring every analytic extension to be formalized at once.

\section{The included rational Hahn reference kernel}\label{sec:prototype}
\subsection{Implemented scope}
The accompanying \code{code/RationalHahn.wl} implements $\Q(t_1,\ldots,t_d)$ with the monomial interpretation \eqref{eq:generators}. It uses ordinary polynomial numerator--denominator expressions with rational coefficients, rejects inexact or transcendental coefficients, and applies the reverse-coordinate exponent order from \cref{prop:core}.

The implemented operations are construction, arithmetic, reciprocal, exact sign and comparison, leading data, valuation, and standard part for finite elements. Domain mismatch and zero division return \code{Failure}. For the valuation of zero, the program uses \code{Infinity} only as the distinguished \emph{valuation marker} $+\infty$, never as the representation of a surreal field element.

The program does not implement the effective real closure, complex pairs, arbitrary exponent groups, lazy infinite supports, transcendental functions, or formal proof checking. Those are subsequent layers of the design, not hidden features of the prototype. Only the public constructor's outputs should be passed to its arithmetic routines; manually forged internal heads are outside its contract.

\subsection{Examples of use}
After loading the package and clearing the formal variables:
\begin{lstlisting}
Get["code/RationalHahn.wl"];
Clear[t, s];
(* t = omega^-1; s = omega^(-omega) *)
a = RHMake[1 + t, t, {t, s}];
b = RHMake[s, t^10000, {t, s}];
RHSign[a]                 (* 1 *)
RHValuation[a]            (* {-1, 0} *)
RHStandardPart[b]         (* 0 *)
RHCompare[b, RHMake[1, 1, {t, s}]]   (* -1 *)
RHExpression[RHInverse[a]]           (* t/(1+t) *)
\end{lstlisting}
The $10000$-power comparison illustrates an exact higher-rank order calculation. Its successful execution is not the proof for every $n$; that universal statement follows from the exponent-group definition and \cref{prop:core}.

The essential leading-polynomial routine is short:
\begin{lstlisting}
leadingPolynomial[p_, variables_List] := Module[{r},
  r = First[SortBy[CoefficientRules[p, variables],
                   Reverse[First[#]] &]];
  {First[r], Last[r]}
];
\end{lstlisting}
The actual code collects and validates polynomials before reaching this routine and handles zero separately. The quotient's leading exponent and coefficient are then obtained by subtraction and division, as in \eqref{eq:leadingrational}. This is the key reason the small kernel can be exact without expanding a single infinite reciprocal series.

\subsection{Verification performed and its limits}\label{sec:verification}
The package's arithmetic definitions and the accompanying 34-check suite were evaluated through the Wolfram Language connector. The reported kernel was version 15.0.1 for Linux x86, dated July 2, 2026. The first evaluation passed 31 of 34 checks. Two failed checks needed \code{Together}, rather than only \code{Cancel}, to combine a sum of rational terms; the endpoint inverse-cosine coefficient check needed its positive square-root coordinate. The three corrected checks all passed in a subsequent kernel evaluation. The distributed suite contains the corrections. The saved report records these evaluations rather than describing them as one untouched all-passing run.

The tests cover 22 rational-kernel behaviors, including invalid-input cases, and 12 finite algebraic or series examples. They do not verify real-closed-field algorithms, infinite support validity, the global phase classification, or transfinite recursions. In particular, a finite series check is evidence about those coefficients only. The support and branch arguments in the article supply the mathematical justification beyond the tested orders.

The connector evaluated the definitions as code; the distribution additionally wraps them in a normal Wolfram package file with usage messages and comments. The arithmetic content is the same, but loading that local path in a separately installed Wolfram system is a user-side reproducibility step, not a claim about an unperformed local installation test. The suite can be run with the supplied \code{RunChecks.wl} driver.

\section{Worked exact calculations across the layers}\label{sec:examples}
The examples below make the difference between an exact descriptor and its finite series view explicit. Let $t=\omega^{-1}>0$. Whenever a finite power series is displayed with $O(t^m)$ in this section, the exact object is specified separately; the omitted terms are not silently zero.

\subsection{A reciprocal, a square root, and exact cancellation}
The exact rational object
\[
 A=(\omega+1)^{-1}=\frac{t}{1+t}
\]
has the view $t-t^2+t^3-t^4+O(t^5)$. Its error after $N$ terms has exact valuation $N+1$. A positive algebraic-root object gives
\[
 B=\sqrt{\omega^2+1}
  =\omega+\frac1{2\omega}-\frac1{8\omega^3}
       +\frac1{16\omega^5}+O(\omega^{-7}).
\]
In contrast to a generic infinite-value system, the rational kernel immediately proves
\[
 (\omega+1)^2-\omega^2-2\omega=1.
\]
No limiting argument, rounding, or series expansion is needed.

\subsection{Finite sine and an infinite slope}
An exact germ descriptor at $\pi/6$ produces
\[
 \sin(\pi/6+t)
 =\frac12+\frac{\sqrt3}{2}t-\frac14t^2
         -\frac{\sqrt3}{12}t^3+O(t^4).
\]
Its standard part is $1/2$. The coefficient field must contain $\sqrt3$, but the series support is simply contained in $\N$.

The exact inverse-tangent reduction is
\[
 \arctan\omega=\frac\pi2-\arctan t
 =\frac\pi2-t+\frac{t^3}{3}-\frac{t^5}{5}
       +\frac{t^7}{7}+O(t^9).
\]
The infinite argument causes no infinite phase choice: inverse tangent returns a finite direction angle. This is different from asking for $\sin\omega$.

\subsection{A multiple analytic zero and a finite polynomial problem}
The coherent equation $\sin^2 z=t$ has exactly two zeros in the monad of zero, counted with multiplicity. They are
\[
 z_\pm=\pm\arcsin(t^{1/2})
 =\pm\left(t^{1/2}+\frac{t^{3/2}}6+
           \frac{3t^{5/2}}{40}+\frac{5t^{7/2}}{112}+\cdots\right).
\]
This is the source analysis manuscript's branch-splitting example. Locally it can be represented as two analytic-root descriptors after ramification. On the full canonical strip modulo $2\pi$, put $U=E_{\mathscr H}(iz)$. The equation becomes
\[
 -U^4+(2-4t)U^2-1=0,
\]
which has four roots counted with multiplicity. The two other angular roots lie near $\pi$ modulo $2\pi$. A local monad solver and a full-strip periodic solver are answering different, compatible questions.

\subsection{An infinitesimal perturbation of a logarithmic cut}
For $z=-1-it$,
\[
 |z|=\sqrt{1+t^2},\qquad
 \LogNo|z|=\tfrac12\Logi(1+t^2).
\]
The anchored logarithm has imaginary part $\pi+\arctan t$; the actual principal logarithm has imaginary part $-\pi+\arctan t$. Their real parts agree and their imaginary parts differ by $2\pi$. Both exponential images under \eqref{eq:stripexp} equal $z$. A branch-keyed cache can store both without inconsistency.

\subsection{A valid nonconvergent ordinary series}
The exact formal-germ value
\[
 F(t)=\sum_{n\ge0}n!t^n
\]
is meaningful by Hahn summability. It satisfies the coefficient identity
\begin{equation}\label{eq:factorialode}
 t^2D_tF+(t-1)F=-1.
\end{equation}
Indeed, for $n\ge1$ the coefficient is $(n-1)(n-1)!+(n-1)!-n!=0$, and the constant term is $-1$. This identity provides an efficient recursive coefficient descriptor.

The homogeneous equation associated with \eqref{eq:factorialode} has the solution $e^{-1/t}/t$ on positive ordinary $t$, as direct differentiation verifies. It is smaller than every power of $t$ there. Thus an asymptotic power-series solution cannot distinguish adding a multiple of this homogeneous term. In a Gonshor-compatible infinite-scale interpretation at $t=\omega^{-1}$, the corresponding positive term is $\omega\ExpNo(-\omega)$, not zero.

It still does not specify a unique ordinary analytic function at positive real $t$ by convergence of that series. Moreover, if a scalar derivation with $\partial t=-t^2$ is chosen, the displayed equation is not obtained by blindly replacing $D_t$ by $\partial$. This example simultaneously tests summation semantics and derivative semantics.

\section{Capability matrix and implementation sequence}\label{sec:roadmap}
\subsection{What can be promised at each stage}
In the table, ``exact node'' means a uniquely interpreted finite descriptor, not a promise of a complete simplifier. ``Effective'' refers to the specified input representation and supplied certificates.
\begin{center}\small
\begin{longtable}{@{}L{0.23\textwidth}L{0.32\textwidth}L{0.36\textwidth}@{}}
\toprule
Layer or operation & Exact representation available & Algorithmic promise and restriction\\\midrule\endhead
Rational monomial core & Polynomial fractions with exact coefficients & Terminating arithmetic, equality, sign, valuation; implemented here over $\Q$.\\[0.35em]
Real closure / complex roots & Finite equations and ordered root selectors & Implementable terminating algebraic operations; not implemented in the reference kernel.\\[0.35em]
Infinitesimal reciprocal and powers & Rational/root node or positive-support series & Coefficients effective with finite dependency procedures; exponent/coefficient promotion allowed.\\[0.35em]
Ordinary analytic lifts & Germ descriptor at $c+\eta$ & Exact under branch and regular-center hypotheses; coefficient procedures may leave ordinary constants symbolic.\\[0.35em]
Real exp and log & Gonshor-semantic graph & All real inputs for exp, positive for log; normalization only where the backend supports the scales.\\[0.35em]
Canonical complex exp & Strip object on $\mathscr H$ & Imaginary part finite; real growth may be infinite.\\[0.35em]
Canonical sine and cosine & Finite-angle / $\mathscr V$ strip object & Real part finite; arbitrary imaginary growth allowed.\\[0.35em]
Infinite circular phases & Named character extension & Exact only after its coherent phase provider is fixed.\\[0.35em]
Inverse real trigonometry & Finite-angle or ramified germ data & Natural real-surreal domains; endpoint branches explicitly selected.\\[0.35em]
Polynomial / trig solving & Algebraic or finite cluster objects & Complete over the effective algebraic core; broader coefficients require effective zero and root capabilities.\\[0.35em]
Function differentiation & Germ/coherent derivative descriptor & Scalars fixed; not a scalar-field derivation.\\[0.35em]
Scalar derivation & Named differential-field structure & Only supported identities and summable families; different from formal $D_t$.\\[0.35em]
Coherent contour integration & Coefficientwise integral descriptor & Common support, ordinary contour, and analytic-domain conditions retained.\\[0.35em]
General lazy Hahn equality & Finite program plus semantic certificate & No total universal procedure, even for $\{0,1\}$ coefficients on $\N$.\\[0.35em]
Numerical display & Coefficient approximations or explicit scale specialization & Not an ordinary numeric value of an infinite surreal.\\
\bottomrule
\end{longtable}
\end{center}

\subsection{A staged implementation plan}
\paragraph{Stage 1: exact ordered rational algebra.}
Implement validated coefficient and exponent domains, rational fractions, arithmetic, equality, sign, valuation, finite standard parts, and serialization. Make the public data contract precise and test its errors as thoroughly as its successful identities. Add surcomplex pairs and exact squared norms. This stage already supports many all-scale geometric calculations.

\paragraph{Stage 2: algebraic extensions.}
Add real-root selectors, ordered extension towers, exact comparison, and complex polynomial roots. Implement finite-dimensional linear algebra and basic polynomial tools over these coefficients. Keep root selection independent of ordinary numerical approximation. Use exact Newton data as a view of a root, not as a substitute for its selector.

\paragraph{Stage 3: support-certified series and jets.}
Add positive-grid constructors, finite coefficient-dependency procedures, reciprocal and binomial expansions, precision propagation, and compressed cutoff objects. Initially restrict support grammars enough that every advertised coefficient request terminates. Preserve exact origin graphs behind truncations.

\paragraph{Stage 4: finite analytic and geometric functions.}
Implement ordinary-center Taylor lifts, finite-angle trigonometry, inverse functions with endpoint ramification, phase classes, logarithm branch policies, and finite trigonometric polynomial algebraization. Reuse host coefficient functions only with domain checks. Add coherent ordinary-contour functionals for a controlled coefficient-function grammar.

\paragraph{Stage 5: exponential and transserial scales.}
Develop the exp--log hierarchy and its compatible surreal interpretation. Add admitted forms of transserial composition and differential operations, each with its own support and leading-term algorithms. Keep general semantic exp--log nodes available even when normalization is incomplete. Treat resurgent special-function extensions and hyperseries as separate advanced providers, not requirements for the earlier stages.

\subsection{Complexity and engineering priorities}
The largest practical risks are expression swell, cancellation-driven searches for leading terms, and uncontrolled domain promotion. A compact rational expression can have a very large series expansion; an algebraic root can have complicated nested coefficients. Store shared exact graphs, cache verified leading data, and make expansion demand-driven.

A total decision procedure need not be the first procedure tried. Fast domain-specific simplifiers can precede an expensive complete algebraic fallback. Outside a decidable domain, every budget exit must remain distinguishable from mathematical falsehood. Persistent caches must include branch, phase, support, and derivation identifiers; reusing a result under a changed interpretation can be an exactness bug even when all printed numbers look identical.

Resource accounting should include coefficient bit sizes, polynomial degrees, extension degrees, exponent-tree sizes, cached coefficient counts, and support-dependency counts. A single ``precision'' integer is inadequate for a hierarchy of infinitesimal scales. Interactive output should show the leading scales and permit inspection of the exact underlying descriptor.

\subsection{Tests that guard semantics, not only arithmetic}
A regression suite should include cancellation to an exact finite value; zeros hidden behind long coefficient prefixes; a mathematically valid but nonenumerable finite-cutoff request; division by exact zero versus unknown zero; standard part of an infinite input; incompatible exponent domains; root branches separated infinitesimally; actual versus anchored logarithms; unsupported infinite circular phases; and differentiation under each named convention.

Tests should also check that \code{Normal} or a similar jet-to-expression conversion cannot silently promote a truncation to an exact equality. Negative tests are essential: rejecting an inadmissible sum or returning a phase obligation is often the correct mathematical answer, not missing functionality.

\section{Conclusions}\label{sec:conclusions}
A substantial surreal CAS is implementable. Its reliable center is an effective ordered field of finitely described monomial expressions, expanded by ordered algebraic extensions and complexification. This already gives exact arithmetic, comparisons, polynomial roots, linear algebra, and a large part of finite geometry at infinite and infinitesimal scales.

Beyond that center, exact representation remains possible without complete normalization. Certified Hahn constructors support genuine infinite sums and local analytic values. Exact germ and implicit-root descriptors support elementary and special functions on controlled domains. A Gonshor-compatible exponential--logarithmic or transseries backend adds infinite growth scales. The supplied manuscripts' finite-angle and canonical strip theory provide a particularly useful trigonometric layer without any need to select $\sin\omega$.

The boundaries are equally concrete. Finite syntax cannot encode every surreal. Arbitrary computable coefficient streams cannot have a universal terminating equality test. Higher-rank cutoff outputs need not be finite lists. Ordinary asymptotic data may omit nonzero exponentially small surreal information. Infinite circular phases, global logarithmic branches, and scalar derivations require explicit structures. None of these boundaries makes the implementable subsets uninteresting; they determine a sound software architecture.

The most useful meaning of ``exact surreal support'' is therefore not a single unchecked promise. It is a visible contract: an exact denotation, a known representation class, a declared set of effective operations, certified remainder information when approximating, and explicit analytic conventions. With those contracts, a symbolic CAS can preserve the genuinely non-Archimedean information that ordinary numeric systems discard while remaining mathematically interpretable and practically extensible.

\appendix
\section{Source-to-design audit}\label{app:audit}
References to theorem names below are to the supplied source files, not claims that the source manuscripts implemented those algorithms.
\begin{center}\small
\begin{longtable}{@{}L{0.27\textwidth}L{0.35\textwidth}L{0.29\textwidth}@{}}
\toprule
Source result or convention & Role in this article & What is not inferred\\\midrule\endhead
Analysis: ``Strong summability'' and Hahn--Neumann lemma & Semantic certificates for products and infinitesimal evaluation & No automatic zero test or effective support decomposer.\\[0.35em]
Analysis: ``Evaluation of arbitrary formal series'' & Existence of sufficiently small analytic scales & No algorithm for extracting a global bound from an arbitrary program.\\[0.35em]
Analysis: ``Hahn-coherent data'' and composition theorem & Common ordinary domains and support-aware function objects & No gluing of arbitrary unrelated fine-local germs.\\[0.35em]
Analysis: preparation and exact root lifting & Implicit roots, finite clusters, coefficient recursions & No completion of arbitrary transfinite recursion by a finite loop.\\[0.35em]
Analysis: contour and moving-pole residue theorems & Coefficientwise exact integration and residue example & No unrestricted fine-topological Riemann integration.\\[0.35em]
Analysis: global logarithm and all-scale rigidity & Separate branch policies and carefully qualified global function types & No universal nonexistence of transcendental functions.\\[0.35em]
Trigonometry: finite-angle circle and inverse functions & Canonical geometry at arbitrary lengths & No necessity for infinite circular angles.\\[0.35em]
Trigonometry: canonical strip functions & Domain-specific complex exp, sine, and cosine & No uniquely specified evaluation beyond the strips.\\[0.35em]
Trigonometry: character classification and infinite periods & Named global phase providers & No default value of $\sin\omega$ forced by local axioms.\\[0.35em]
Present article: \cref{prop:countability,thm:halting,prop:core,thm:cutoff} & Explicit computability and storage boundaries; a decidable core & No priority claim for these elementary deductions.\\
\bottomrule
\end{longtable}
\end{center}

\section{Minimal contracts for a backend}\label{app:contracts}
A coefficient backend should specify exact construction, zero and sign capabilities, arithmetic, conversion to larger coefficient domains, and whether its ordinary constants admit certified numerical enclosures. A monomial backend should specify equality, multiplication, order, rational-root promotion, and the embedding used to interpret it inside the surreal field.

A support backend should distinguish a semantic proof of well ordering from algorithms for membership, finite coefficient dependencies, leading-term discovery, and cutoff materialization. It should declare which of these operations are total. An analytic backend should specify the ordinary or surreal domain, branch and extension identifiers, Taylor coefficients, derivative mode, and admissible composition rules.

An operation result should preserve its exact origin and obligations. If the answer is a truncation, it must include a remainder descriptor and the error calculus used to obtain it. If the answer is conditional, the condition must remain visible in further transformations. If an unsupported request has a well-defined semantic node, the system may retain that node without pretending to have evaluated or normalized it.

\section{Reproducibility and archive contents}\label{app:archive}
The archive contains this self-contained LaTeX source and its compiled PDF; the rational Hahn package, its corrected 34-check suite, and a driver; a report of the actual Wolfram connector evaluations; and the two supplied source manuscripts in a separate \code{sources} directory for provenance. Those source manuscripts are included unchanged. Their own referenced verification scripts or additional archives were not supplied and are not represented as included or rerun here.

Compile this article with a standard PDFLaTeX installation and the listed common packages, for example by running \code{pdflatex surreal_cas.tex} twice. No bibliography processor or external graphics are needed. The code driver requires Wolfram Language and can be run from the extracted archive. Its coefficient and example calculations are exact finite tests; a successful run does not upgrade them into proofs of the infinite constructions.

\begingroup\small\raggedright
\begin{thebibliography}{99}
\addcontentsline{toc}{section}{References}
\setlength{\itemsep}{0.4em}
\bibitem{SourceAnalysis}
\emph{Surcomplex Analysis: Infinitesimal Calculus, Hahn-Coherent Holomorphy, and Contour Theory}.
Supplied manuscript, \code{surcomplex_analysis(2).tex}, September 21, 2026.
Primary basis for the analysis conventions used here. Source labels include
\code{def:strong}, \code{thm:smallscale}, \code{def:coherent},
\code{thm:preparation}, \code{thm:globallog}, and \code{thm:rigidity}.
\bibitem{SourceTrig}
\emph{Trigonometry on the Surcomplex Plane: Canonical Angles, Infinite Triangles, Infinitesimal Contact, and Hahn-Analytic Oscillation}.
Supplied manuscript, \code{surcomplex_trigonometry.tex}, September 21, 2026.
Primary basis for finite-angle, strip, branch, root-count, and infinite-phase conventions.
Source labels include \code{thm:polar}, \code{thm:tangency}, \code{thm:stripexp},
\code{thm:polyroots}, \code{thm:characters}, and \code{thm:infiniteperiods}.
\bibitem{Conway}
John H. Conway. \emph{On Numbers and Games}. Academic Press, 1976; second edition, A K Peters, 2001. Foundational reference cited by both supplied manuscripts.
\bibitem{Gonshor}
Harry Gonshor. \emph{An Introduction to the Theory of Surreal Numbers}.
London Mathematical Society Lecture Note Series 110. Cambridge University Press, 1986.
\href{https://doi.org/10.1017/CBO9780511629143}{doi:10.1017/CBO9780511629143}.
\bibitem{vdDE}
Lou van den Dries and Philip Ehrlich. ``Fields of surreal numbers and exponentiation.''
\emph{Fundamenta Mathematicae} 167 (2001), 173--188.
\href{https://doi.org/10.4064/fm167-2-3}{doi:10.4064/fm167-2-3}.
Erratum, 168 (2001), 295--297. This article uses the normal-form and exponential-structure background, not unqualified birthday estimates from the original paper.
\bibitem{Poonen}
Bjorn Poonen. ``Maximally complete fields.''
\emph{L'Enseignement Math\'ematique} (2) 39 (1993), 87--106.
\href{https://math.mit.edu/~poonen/papers/amsval.pdf}{Author's full text}, especially Corollary 4.
\bibitem{PerrucciRoy}
Daniel Perrucci and Marie-Fran\c{c}oise Roy. ``Elementary recursive quantifier elimination based on Thom encoding and sign determination.''
\emph{Annals of Pure and Applied Logic} 168 (2017), 1588--1604.
\href{https://arxiv.org/abs/1609.02879}{arXiv:1609.02879v2};
\href{https://doi.org/10.1016/j.apal.2017.03.001}{doi:10.1016/j.apal.2017.03.001}.
\bibitem{vdH}
Joris van der Hoeven. \emph{Transseries and Real Differential Algebra}.
Lecture Notes in Mathematics 1888. Springer, 2006.
\href{https://www.texmacs.org/joris/ln/ln.html}{Author's online text}, especially the chapters on grid-based series, Newton polygons, and operations on transseries.
\bibitem{BMDerivations}
Alessandro Berarducci and Vincenzo Mantova. ``Surreal numbers, derivations and transseries.''
\emph{Journal of the European Mathematical Society} 20 (2018), 339--390.
\href{https://doi.org/10.4171/JEMS/769}{doi:10.4171/JEMS/769};
\href{https://arxiv.org/abs/1503.00315}{arXiv:1503.00315}.
\bibitem{BMFunctions}
Alessandro Berarducci and Vincenzo Mantova. ``Transseries as germs of surreal functions.''
\emph{Transactions of the American Mathematical Society} 371 (2019), 3549--3592.
\href{https://doi.org/10.1090/tran/7428}{doi:10.1090/tran/7428};
\href{https://arxiv.org/abs/1703.01995}{arXiv:1703.01995v3}.
\bibitem{CostinEhrlich}
Ovidiu Costin and Philip Ehrlich. ``Integration on the Surreals.''
\emph{Advances in Mathematics} 452 (2024), article 109823.
\href{https://arxiv.org/abs/2208.14331}{arXiv:2208.14331v5}.
A distinct constructive extension/integration framework, not identified here with the supplied manuscript's coefficientwise contour functional.
\bibitem{Hyperseries}
Vincent Bagayoko and Joris van der Hoeven. ``Surreal numbers as hyperseries.''
\href{https://arxiv.org/abs/2310.14879}{arXiv:2310.14879v1}, October 23, 2023.
Cited for its mathematical representation theorem, not for a finite universal encoding or implementation claim.
\bibitem{Mantova2026}
Vincenzo Mantova. ``Monotonicity and a Taylor approximation theorem for transseries.''
\href{https://arxiv.org/abs/2601.07747}{arXiv:2601.07747v2}, revised May 9, 2026.
Preprint record consulted September 21, 2026.
\bibitem{WLSeries}
Wolfram Research. ``SeriesData.'' \emph{Wolfram Language Documentation}.
\url{https://reference.wolfram.com/language/ref/SeriesData.html}.
\bibitem{WLRoot}
Wolfram Research. ``Root.'' \emph{Wolfram Language Documentation}.
\url{https://reference.wolfram.com/language/ref/Root.html}.
The consulted documentation includes general-equation neighborhood representations as well as polynomial indexing.
\bibitem{WLSpecial}
Wolfram Research. ``Mathematical Functions: Working with Special Functions.''
\emph{Wolfram Language Documentation}.
\url{https://reference.wolfram.com/language/tutorial/MathematicalFunctions}.
\bibitem{WLInactive}
Wolfram Research. ``Inactive.'' \emph{Wolfram Language Documentation}.
\url{https://reference.wolfram.com/language/ref/Inactive.html}.
\bibitem{WLUpValues}
Wolfram Research. ``UpValues.'' \emph{Wolfram Language Documentation}.
\url{https://reference.wolfram.com/language/ref/UpValues.html}.
\bibitem{WLNumericQ}
Wolfram Research. ``NumericQ.'' \emph{Wolfram Language Documentation}.
\url{https://reference.wolfram.com/language/ref/NumericQ.html}.
\bibitem{SageLazy}
The Sage Developers. ``Lazy Series.'' \emph{SageMath Reference Manual}.
\url{https://doc.sagemath.org/html/en/reference/power_series/sage/rings/lazy_series.html}.
\bibitem{SageAsymptotic}
The Sage Developers. ``Asymptotic Ring.'' \emph{SageMath Reference Manual}.
\url{https://doc.sagemath.org/html/en/reference/asymptotic/sage/rings/asymptotic/asymptotic_ring.html}.
\end{thebibliography}
\endgroup
\end{document}
